
%the propensiy, the issue of determinis by that of realis, Stull Popper appears never to have abandoned the conviction that the very notion of past trajectories was not merely a \s{matter of taste}, nor was it meaningless, as Schlick maintained. They are an essential component of \qm. For Popper, Heisenberg's concession was the crack that  \s{Copenhagen interpretation} measuring agencies with the physical object under investigation


%
%Hermann’s long study appeared in the \emph{Abhandlungen der Fries’schen Schule} and was also distributed as a separately printed booklet or offprint.

%the momentum (position) operator

%The statistical interpretation therefore does not consist in assigning a wave-like probability distribution to an underlying ensemble of particles. Rather, it becomes applicable only when one considers the probabilistic relations governing the transition from one observational context to another.




%probability is binary (relational)—it links two contexts, the preparation context \psi and the measurement context \phi.


%\citep{Hermann1935b}

%$(\psi, \phi) \rightsquigarrow P(\phi \mid \psi)$.

%($\psi \rightsquigarrow P(p, q)$)

%gives

%The statistical character of quantum mechanics therefore does not reside in the fact that the wave function only refers to an ensemble of systems but in the probabilistic relation between complementary experimental contexts.


%In Popper's interpretation, the wave-description refers to the probability distribution over an ensemble, while the particle description is an individual member of that ensemble.   For Hermann the wave and particle descriptions concern the same individual system limiting the the contextual natura Consequently, the wave function does not describe a probability distribution over particles. It assigns the probabilities of the outcomes that would be obtained in an incompatible experimental arrangement.


%For Popper the wave function $\varphi$ a wave-like probability distribution over an ensemble \origg{Schar} of \emph{particles} with unknown position and momenta \citep[\S5]{Hermann1935}. For this reason acoording. Hermann objects that this misconstrues the significance of wave-particle duality applies to every single system.

%For this reason Popper: The statistical interpretation therefore applies not within a single context but only to the transition between incompatible contexts. A state prepared in the momentum context determines only the probabilities of the outcomes in a position context, and conversely:


%Acoording Hermann objects that this misconstrues the significance of wave-particle duality. Wave and particle are not alternative descriptions of different members of an ensemble, but complementary descriptions of the same individual system in different experimental contexts.

%Popper interprets the wave function as a probability distribution over an ensemble of particles existing within a single experimental arrangement.

%statistical interpretation of \qm rests on the idea that the function $\varphi$ describes a wave-like probability distribution over an ensemble \origg{Schar} of \emph{particles} \citep[\S5]{Hermann1935}. According to Hermann, however, this interpretation misconstrues the significance of wave-particle duality. The wave and particle descriptions apply to the same individual physical system, although which description is appropriate depends on the experimental context:

%Like \citep{Schlick1931}.%However, Hermann's contextualist reading is more subtle than Schlick's blunt operationalist approach. For Schlick, a past trajectory cannot be tested; for Hermann, the two context-dependent determinations cannot be combined
%


%For Hermann, wave–particle duality has two distinct consequences. At the inter-contextual level, it determines which measurement can be performed: one must choose an experimental arrangement in which the object is described either under the conditions appropriate to a position measurement or under those appropriate to a momentum measurement. The two contexts cannot be fused into a single, context-independent description of the system.
%
%%Hermann contrasts this illegitimate reconstruction with the interpretation of a measurement, the \qt{retrospective inferences}{Rückschlüssen} required for the interpretation of a measurement.

%However, the resolving power of the focal-plane mark is limited to a narrow angular spread $\delta\vartheta$ of the refracted light. This instrumental limit determines the irreducible uncertainty $\delta\alpha$ of the initial scattering angle.

%All light entering the lens passes through both the focal and image planes, but different families of rays converge in each plane at $P'$\ref{fig:HW}. In the arrangement on the left, rays arriving at the lens from the same direction, originating from different points $P_1, P_2, P_3 \ldots$, are focused at a single point $P'$ on the screen $S$, which is positioned at the focal distance $f$. In the arrangement on the right, each point $P$ in the object plane $L$ is mapped onto a corresponding point $P'$ in the image plane $S$, located further away at the image distance $d_i$. The lens brings parallel rays to a focus at the focal distance. Rays emanating from a finite-distance object at $P$ arrive at the lens already diverging; the lens must first compensate for this divergence so that the rays converge only at a greater distance, at the image distance $d_i$, after which they diverge again.




%His reply to point $3$ is strikingly brief: he simply states that he fully agrees. The point that would become central to Hermann’s public refutation of Popper is therefore already present in her letter, while \vW merely endorses it without further elaboration.  He cannot think of a single experiment that would count as proof of the wave picture.


%The decisive point was structural, not merely practical.

% \VW\ stressed that this was the reverse of Heisenberg's example of an initial position determination followed by a momentum measurement in a field. \citep[21--22]{Heisenberg1930} In both temporal directions, the uncertainty relation blocks the combination of sharp momentum knowledge and retrodictive position knowledge required by Popper's argument.

%\subsection{Heisenberg's Countereply}
%
%In his reply of \datedm{10}{12}{1934}, Heisenberg enclosed \vW's criticism and opened by assuring Popper that, after discussing the matter with \vW, he had had the note corrected so that it no longer appeared unjustifiably sharp.\letterKPAp{Heisenberg}{Popper}{10}{12}{1934}[305-32] Heisenberg summarized \vW's result as follows: the proposed \s{Aussonderung} of electrons by an electric field is not a genuine momentum measurement in Popper's sense, since the experiment cannot yield an exact trajectory of the particle before the momentum measurement. He then turned to Popper's analogous suggestion of using a light filter instead of an electron spectrometer. As he put it, the optical case could be treated in the same way. The filter would select light of a given frequency, but it would not prepare quanta with an infinitely sharp momentum.
%
%Heisenberg suggested that Popper consider reflective resonance filters, such as sodium or mercury vapor, following Wood's \origg{Reststrahlenverfahren}. Such a device reflects or re-emits radiation near the resonance frequency, while other components of the incident light are not selected in the same way. Thus a broad beam can be reduced, for example, to the sodium D-lines. Yet the selected light is never strictly monochromatic. Since the filtering process depends on the excitation and subsequent de-excitation of the reflecting atom, the resonance line has a finite width. Heisenberg expresses this by writing that the inaccuracy of the momentum of the light quantum is given by
%\begin{equation*}
%\Delta p = \frac{h\Delta\nu}{c^2}.
%\end{equation*}
%The precise form of the factor is less important for his argument than the point that a finite frequency width $\Delta\nu$ entails a finite momentum width $\Delta p$.
%
%Heisenberg then connects this linewidth to a temporal indeterminacy. The probable residence time of the light quantum in the reflecting atom is determined by the lifetime of the corresponding excited state:
%\begin{equation*}
%\Delta t = \frac{1}{\Delta\nu}.
%\end{equation*}
%Thus, even if the later position of the light quantum were known with precision, its position before reflection could not be reconstructed exactly, because the moment at which it was in the reflecting atom remains indeterminate by $\Delta t$. This temporal indeterminacy produces an uncertainty in the reconstructed position before reflection, and Heisenberg concludes that
%\begin{equation*}
%\Delta x \Delta p \sim h.
%\end{equation*}
%In the optical case, therefore, the uncertainty relation is preserved for the same structural reason as in \vW's analysis of the electrostatic analyzer: the physical procedure that selects momentum also prevents an exact retrodiction of the earlier position.
%
%Heisenberg finally made explicit that Popper's \s{nichtprognostischen Messungen} had not, in his view, been shown to be impossible in general merely by \vW's discussion. Rather, what had been shown was that the particular examples Popper had proposed did not work. He nevertheless thought it plausible that the impossibility of such measurements was a very general feature of quantum mechanics. He remarked that a proof of this point would presumably amount to the content of an important paper in mathematical physics, and added that he feared Popper had here followed one of those misleading errors by which he himself, Bohr, and others had been tempted since 1927, before they came to believe that they had really understood quantum theory. He closed by thanking Popper for sending his book, which he had not yet been able to study carefully.\letterKPAp{Heisenberg}{Popper}{10}{12}{1934}[305-32]


% As he put it, the optical case could be treated in the same way. The filter would select light of a given frequency, but it would not prepare quanta with an infinitely sharp momentum, the time of the passage wouild be lost is the time of pasage, since the filter ....  In his reply of \datedm{16}{12}{1934},



%Otherwise an electron arriving higher on the screen might simply have entered the field at a different position rather than having had a different momentum.


%between the \s{non-predictive measurements} he has defined and testable measure about the past.

%Popper's \s{non-predictive measurements} evade the uncertainty relations only because the statements reporting their results do not describe physically possible measurements. Verifiable inferences about the past, by contrast, remain subject to the same accuracy limits as verifiable predictions about the future, given the time-reversal symmetry of the laws of quantum mechanics.


%

%\q{You will then see what kind of criticism it is. Essentially, it concerns the concept of the \s{non-predictive measurement,} which in our view appears to be misunderstood in your work}. Heisenberg expressed confidence that an agreement could be reached.

%Popper, who at the time had only one short publication to his name, was still unknown, and it is therefore quite remarkable that Heisenberg chose to engage with him, enlisting \vW for the purpose---the most philosophically inclined of his collaborators.

%The experiment described by Mr. Popper, when discussed correctly, does not permit one to fall below the limit of accuracy given in the uncertainty relation. It is not possible to determine precisely from the measurements in X the place and time of the collision in S. To see this, however, it is necessary to provide precise information about the nature of the measurements in X. Suppose, for example, that first the momentum of the A-particle is measured and then the time of its arrival at a given place. The position measurement does indeed destroy the knowledge of the momentum after the collision; by contrast, from the momentum before the collision and the position, a “trajectory” of the particle before the collision can be constructed. This fact is presumably what Mr. Popper intends to express by the concept of a “non-prognostic measurement.” This trajectory, however, is in principle uncontrollable. For it holds only for the time interval between the end of the momentum measurement and the beginning of the position measurement, during which the particle has no interaction at all with its environment, and it cannot be extended back into the period before the momentum measurement, since the latter in turn destroys the knowledge of the position in accordance with the uncertainty relation.

%The position measurement does indeed destroy the knowledge of the momentum after the collision; by contrast, from the momentum before the collision and the position, a \q{trajectory} \origg{Bahn} of the particle before the collision can be constructed. \q{This trajectory, however, is in principle uncontrollable} \citep{Weizsäcker1934}.



%The problem of trajectory reconstruction turned out to be a special case of the problem of state attribution.

%Hermann’s critique of Popper may be understood as the temporal extension of her analysis of state attribution. Her original point is that the attribution of a complete classical state at a given instant is context-dependent. But a trajectory is nothing over and above a temporally ordered sequence of such state attributions. Consequently, if each individual state attribution depends on the experimental context, then no context-independent trajectory can be reconstructed by stringing them together.


%\q{The uncertainty principle does not forbid calculating momentum from past data}.



%See Bohr (1985, pp. 415–418), for the German original and pp.  for the English translation

%Bohr, N. (1985). Collected works. In J. Kalckar (Ed.), Foundations of quantum physics i (Vol. 6) (pp. 1926–1932). Amsterdam: North-Holland.



\todo{Feyerabend}
% \hide{that, in a one-slit experiment, a definite trajectory can be reconstructed retrospectively from the position preparation and the subsequent momentum measurement. In modern preparation--measurement language, Hermann would object that the position preparation determines a diffracted wave, not a particle with an unknown momentum component. The subsequent measurement licenses the attribution of a definite particle momentum to the outcome, but not the retrospective attribution of a trajectory}


%After the slit the system is a spherical wave, in a position preparation context there is no determinate momentum. After the The context only allows to interpret the spot on the screen as a momentum measurement, rather than a position measurement

%the preparation merely reveals a pre-existing trajectory.
%the preparation defines the context within which a causal reconstruction is meaningful. If Hermann’s “context” is interpreted as what we now call a preparation procedure, then the contextuality lies at the source, not primarily at the detector. The preparation fixes the state (or ensemble); the subsequent measurement merely samples one of the complementary observables. That seems closer to Hermann’s actual discussion than the modern tendency to associate contextuality exclusively with the measurement apparatus.

%With a fixed slit, the screen coordinate can be interpreted as a momentum measurement because the initial position is well defined. With a movable slit, the initial position is no longer well defined, so the same screen coordinate loses its unambiguous momentum interpretation and is naturally regarded simply as a position detection.

%Consider a single slit experiment:


%    relative to the actual preparation and measurement performed


%The single slit prepares the particle in a position context. The screen then performs a momentum measurement. From the initial position preparation and the final momentum outcome one may reconstruct the particle’s intermediate trajectory. However, this reconstruction is valid only within that particular experimental context. Had the final measurement instead been one of position, a different reconstruction would have been obtained. Thus what is reconstructed is not a context-independent history of the particle but the history associated with the actual preparation-and-measurement sequence.


Using the notation of the period, one may write the post-collision state as

\begin{equation*}
\psi_{EP}(x_E,x_P)
=
\sum_{mn} c_{mn}\psi^E_m(x_E)\psi^P_n(x_P).
\end{equation*}

This state is factorizable only in the special case in which the coefficients satisfy $c_{mn}=a_m b_n$. In general, this condition is not satisfied. The collision therefore produces an entangled state in which the electron and photon can no longer be assigned independent quantum states. In the idealized case of perfect one-to-one correlations, the state takes the simpler form

\begin{equation*}
\psi_{EP}(x_E,x_P)
=
\sum_n c_n\psi^E_{m(n)}(x_E)\psi^P_n(x_P).
\end{equation*}

The coefficients determine the probabilities of the correlated outcomes: $\left|c_n\right|^2$ is the probability that the photon is found in the state $\psi^P_n$ and the electron in the uniquely correlated state $\psi^E_{m(n)}$. 


Using the notation of the period, one may write the post-collision state as

\begin{equation*}
\psi_{EP}(x_E,x_P)
=
\sum_{mn} c_{mn}\psi^E_m(x_E)\psi^P_n(x_P).
\end{equation*}

This state is factorizable only in the special case in which the coefficients satisfy $c_{mn}=a_m b_n$. In general, this condition is not satisfied. The collision therefore produces an entangled state in which the electron and photon can no longer be assigned independent quantum states. In the idealized case of one-to-one correlations, the coefficients $c_{mn}$ reduce to a single-index set $c_n$.  The probability distribution collapses to a single outcome, so measuring the photon determines the state of electron with certainty.


%P(m,n)=
%\begin{cases}
%|c_n|^2,& m=m(n),\\
%0,&\text{otherwise}.
%\end{cases}


%P(m,n)=
%\begin{cases}
%|c_n|^2,& m=m(n),\\
%0,&\text{otherwise}.
%\end{cases}



%The coefficients determine the probabilities of the correlated pairs: \left|c_n\right|^2 is the probability that the photon is found in the state \psi^P_n and the electron in the uniquely correlated state \psi^E_{m(n)}. Consequently, once the photon is found in the state \psi^P_n, the electron is found with certainty in the state \psi^E_{m(n)}.

%c_{mn} determine the joint probability distribution of pairs

%there is only one allowed electron state for each photon state. 

%More generally, without assuming perfect one-to-one correlations, the coefficients $c_{mn}$ determine the joint probabilities of the measurement outcomes. The quantity $\left|c_{mn}\right|^2$ is the probability that the electron is found in the state $\psi^E_m$ and the photon in the state $\psi^P_n$, with
%
%\begin{equation*}
%\sum_m\sum_n \left|c_{mn}\right|^2=1.
%\end{equation*}

%Accordingly, if the photon is found in the state $\psi^P_n$, the conditional probability that the electron is found in the state $\psi^E_m$ is
%
%\begin{equation*}
%P(m\mid n)
%=
%\frac{\left|c_{mn}\right|^2}
%{\sum_m\left|c_{mn}\right|^2}.
%\end{equation*}


%The former escape the \IR because they are not physical measurements at all, insofar as they cannot be used as initial conditions for existing dynamical laws to make testable predictions. The latter, by contrast, do satisfy the \IR, which von Weizsäcker regarded as symmetric for the past and the future.
 
 %Popper's counterargument is indeed based on a modification of the momentum measuring that he already proposed in Appendix~VI of the \LdF. Popper was surprised at Heisenberg and \vW's reaction. He considered the possibility of reconstructing post trajectories as obvious given that \citep{Heisenberg1930}, on p.~15 in his textbook, concedes that the belief in the existence of such trajectories was a \s{matter of taste} and was not forbidden by the \UR.
%
%Popper must have realized early on that the tip counter was not the right measuring apparatus and that \vW\ was right that a better way was to measure the momentum \emph{measurement} at $X$ via the Doppler effect, that is, by using a collision with a light quantum. According to Popper, however, the latter method has the drawback that it \q{disturbs the particle's momentum (recoil),} so that \qt{the exact instant of the disturbance is not known}{den Impuls des Teilchens stört (Rückstoß) und daß, da der genaue Zeitpunkt der Störung nicht bekannt ist} \autocite[406]{Popper1934b}. Indeed, no reconstruction of the trajectory before the scattering event at $X_1$ is possible in this way. However, following the same argument he had developed in Appendix VI, Popper suggests that one can resort to \q{a momentum \emph{selection} (electron spectrometer, light filter),} which \qt{unlike the momentum measurement of individual particles, for example by means of the Doppler effect---has the property of not affecting the momenta or momentum components of the selected particles}{eine Impuls\emph{aussonderung} (Elektronenspektralapparat, Lichtfilter) hat nämlich- im Gegensatz zu der Impulsmessung individueller Teilchen, etwa mit Hilfe des Dopplereffektes - die Eigenschaft, die Impulse bzw. die Impulskomponenten der ausgesonderten Teilchen nicht zu beeinflussen} \autocite[406]{Popper1934b}. If $[A]$ is an electron beam, one can use a spectra analyzer; in this case \qt{only electrons with a certain \myemph{predetermined} amount of momentum enter the Geiger counter}{nur Elektronen mit einem gewissen vorgegebenen Impulsbetrag in den Spitzenzähler einfallen}. 

%If $[A]$ is a light beam, the filter lets through only those light quanta that \emph{already} possessed the required momentum before entering it  \autocite[406]{Popper1934b}.\todo{selfplagiarism}


%\VW\ began by clarifying that he had not intended to disparage Popper or his book, though he maintained his substantive criticism of Popper's notion of a \s{non-prognostic measurement}. In his view, that notion conflicted with the foundations of quantum mechanics as understood by Heisenberg and himself.\letterKPAp{Von Weizsäcker}{Popper}{6}{12}{1934}[360-21]


%Popper treats the statistical formalism as referring to ensembles of particles with retrospectively reconstructible trajectories; Hermann treats it as referring to ensembles of possible outcomes generated under specified measurement conditions.

%
%Popper: frequency predictions about context-independent reconstructed trajectories.
%
%Hermann: frequency predictions about context-dependent reconstructed measurement outcomes.

%The probabilistic interpretation of the wave function can be understood both intra-contextually and inter-contextually. Thus, in a momentum context, the sharper the momentum determination, the broader the probability distribution for position: one can specify only the probability of finding the particle at a given place. 




%he probabilistic use of the wave function, by contrast, has an inter-contextual function\todo{improve}: it specifies what would be found if a different, incompatible context were realized. If a momentum context is given, the wave function gives only probabilities for the positions that would be obtained under a position measurement. Hence there is no context-independent trajectory. One may in principle reconstruct a \s{momentum history}, but not a corresponding \s{position history}, except probabilistically and vice versa. 




%In this case the first momentum selection at $X$ seemed not to produce any disturbance analogous to that involved in a position measurement; it appeared merely to allow particles with a given momentum to pass through the filter without altering that momentum. The very possibility of Popper's \TE rested on this rather weak premise.\footnote{The apparent asymmetry arises from the fact that position is not an eigenstate of the Hamiltonian operator of a free particle, whereas momentum is. Yet the point Popper missed is that a momentum preparation still projects the system into a momentum eigenstate, rendering its position completely undetermined, since position is not an eigenstate of the momentum operator. Consequently, there is no way to combine information from a preparation and a measurement to reconstruct an individual trajectory before the momentum selection, as Popper had hoped.}

%
%Yet the point Popper missed is that a momentum preparation still projects the system into a momentum eigenstate, rendering its position completely undetermined, since position is not an eigenstate of the momentum operator. 



%Popper explained to Heisenberg his distinction between \begin{inparaenum}[(1)] \item the \s{selection} of an ensemble of particles and the measurement of individual particles. \end{inparaenum}. To determine the past of a particle, one can proceed in three possible combinations of selections and measurements: \begin{inparaenum}[(a)] \item two position measurements, a momentum measurement followed by a position measurement, or a position measurement followed by a momentum measurement \end{inparaenum}. Popper insisted that case $b$ was special. It is true that the \s{momentum measurement} suggested by \vW disturbs the momentum. However, Popper suggests that one can use a \s{momentum selection} e.g. by means a spectral analyzed that leaves the electron's momentum unchanged. In particular, Popper sent Heisenberg the galley proofs of Appendices~VI and~VII of the \Lo, in which he explained this arrangement in details. .


%However, following the same argument he had developed in Appendix VI, Popper suggests that one can resort to \q{a momentum \emph{selection} (electron spectrometer, light filter),} which \qt{unlike the momentum measurement of individual particles, for example by means of the Doppler effect---has the property of not affecting the momenta or momentum components of the selected particles.}{eine Impuls\emph{aussonderung} (Elektronenspektralapparat, Lichtfilter) hat nämlich- im Gegensatz zu der Impulsmessung individueller Teilchen, etwa mit Hilfe des Dopplereffektes - die Eigenschaft, die Impulse bzw. die Impulskomponenten der ausgesonderten Teilchen nicht zu beeinflussen} \autocite[406]{Popper1934b}. If $[A]$ is an electron beam, one can use a spectra analyzer; in this case \qt{only electrons with a certain \myemph{predetermined} amount of momentum enter the Geiger counter}{nur Elektronen mit einem gewissen vorgegebenen Impulsbetrag in den Spitzenzähler einfallen}. 

%However, the tables seem to have turned. Grand synthetic philosophies of science have come to appear suspect, and Popper's philosophy, among others, has fallen out of fashion. At the same time, Hermann's interpretation of \qm has acquired new visibility, as one of the most sophisticated and cohere expression of what we call \s{Copenhagen interpretation}.

%\citets{Popper1967} is now often regarded as a rather naive critic of the Copenhagen interpretation of \qm, whereas Hermann appears as one of its most sophisticated representatives.

%was directed against what was, in many respects, a philosophically naive version of that interpretation \citep{Howard2012}. Hermann's work, by contrast, represents one of its most sophisticated formulations: not a simple denial of causality or reality, but a contextual reconstruction of both within the conditions imposed by quantum measurement.

%\citets{Popper1967} critique of the Copenhagen interpretation of \qm as fundamentally naive \citep{Howard2012}; Hermann has its most refined expression of it as the most sophisticated expression of that interpreation.

%Grand synthetic philosophies of science have come to appear suspect, and Popper's philosophy, among others, has fallen out of fashion. At the same time, Hermann's analysis has acquired new visibility, as philosophy of physics has become increasingly technical.

%However, the tables seem to have turned. Whereas \citep{Popper1967} has remed will become the outmost critiques of the Copenhagen interpretation of \qm as fundamentally, \citep{Howard2012}; Hermann has its most refined expression of it\todo{improve}.



%Grand synthetic philosophies of science have come to appear suspect, and Popper's philosophy, among others, has fallen out of fashion. At the same time, Hermann's analysis has acquired new visibility, as philosophy of physics has become increasingly technical. Yet the tables may turn again, as philosophy of the special sciences seems increasingly unable to connect local technical questions to philosophical issues of wider significance.
 
 



%even we cannot predict them

%the particle always has position and momentum; and the theory does not predict them sharply simply because i t is a probabilistic theory 


%%Popper continued to consider Heisenberg's concession revealing, although he did not appreciate their role in the theory: 
%
%%\q{for now there can be no question whether, according to the quantum theory, an electron can \s{have} a precise position and momentum. It can.} \citep{Popper1967}. 



%The two disputes are therefore connected. If EPR is right, then Popper can claim that past trajectory reconstructions are possible: once position and momentum may be attributed to a system independently of the measurement context, there is no principled reason to deny their attribution at different times and hence their combination into a trajectory. If EPR is wrong, however, Bohr and Hermann can reply that such trajectory reconstructions are impossible. Position and momentum are not simply properties that can be detached from their experimental conditions and recombined at will. They belong to mutually exclusive measurement contexts. Thus one cannot combine a position measurement and a momentum measurement into a single description of the same particle, whether at one time, as in EPR, or across time, as in Popper's reconstruction of past trajectories.

%It entails only the prohibition of diachronic trajectory-attribution in the classical sense, where the trajectory presupposes definite position and momentum as elements of one continuous, context-independent history. 


%According to Popper, Bohr changed the game halfway through by refusing this combination. Whatever one might think, that Bohr's move is truly against synchronizing state attributes, this is the same move that Hermann made against diachronical trajectory reconstruction. 

%simultaneous state attribution

% The question but whether two counterfactual measurement contexts may be fused into one state attribution. 



%\q{It becomes meaningless (if I understand it properly) to attribute to $A$ a momentum if it is unrelated to the momentum framework, or a position if it is unrelated to the position framework.)}. Hermann had applied the same reasoning to a succession of measurements, 

% EPR interpret these alternative conditional assignments as revealing properties that $A$ already possesses independently of the measurement performed on $B$. Popper seems to the brutst of Bohr's reply Bohr denies that the two counterfactual attributions can be combined, since they belong to mutually exclusive experimental contexts. Bohr  deny this step. The two assignments belong to different experimental contexts and cannot be combined into a single context-independent description of $B$. Hermann, to trajevories. 

%A measurement of the position of $B$ permits one to assign a corresponding position state to $A$; a measurement of the momentum of $B$ permits one to assign a corresponding momentum state to $A$. Einstein and 



%2. we have to choose \q{between an experimental framework which would allow us to relate $B$ to a position space, and another experimental framework which would allow us to relate $B$ to a momentum space} We cannot combine the two spaces or co-ordinate systems into one. 


%From the fact that the position of $A$ could be inferred if the position of $B$ were measured, and that the momentum of $A$ could be inferred if the momentum of $B$ were measured, EPR conclude that $A$ possesses both position and momentum at one and the same time. Bohr denies that the two counterfactual attributions can be combined, since they belong to mutually exclusive experimental contexts.



%Bohr's objection to EPR state-attribution has the same structure as Hermann's objection to Popper's trajectory reconstruction. In both cases, the realist inference moves from alternative counterfactual measurement possibilities to a single unified description


%In EPR, the move is synchronic: 





%1. If we decide to measure the position of $B$ we can then calculate the position of the far away $A$. If we decide to measure the momentum of $B$, we can then calculate the momentum of the far away $A$. Thus $A$ must have both position and momentum, even though the theory does not enable us to calculate both predictively at the same time. \q{This beautiful and simple argument seems to me conclusive.}.  

%2. we have to choose \q{between an experimental framework which would allow us to relate $B$ to a position space, and another experimental framework which would allow us to relate $B$ to a momentum space} We cannot combine the two spaces or co-ordinate systems into one. \q{It becomes meaningless (if I understand it properly) to attribute to $A$ a momentum if it is unrelated to the momentum framework, or a position if it is unrelated to the position framework.)}. 

%According to Popper Bohr schifted the argumetn \s{is as gratuitous as the assertion}





%Einstein: We may choose arbitraily whether to measure the position of $B$ or its momentum. 

%  This is a change of the theory. 
%
%he measurement context fixes which conditional description is legitimate, and excludes the counterfactual combination with the incompatible one.
%
%The two antecedents describe mutually exclusive experimental situations. The fact that each attribution is licensed in its own context does not entail that both attributions may be combined into one context-independent description of A.
%
%
%the EPR argument, in Popper’s formulation, depends on moving from two counterfactual possibilities to one simultaneous state attribution.
%
%
%His objection is that Bohr’s appeal to mutually exclusive frameworks misses the point, because Einstein’s argument presupposes precisely this freedom of alternative choice.
%
%Rather, she would say that the causal reconstruction is possible only within the actually chosen measurement context. 

%The analogy may be stated as follows. Bohr's objection to EPR state-attribution has the same structure as Hermann's objection to Popper's trajectory reconstruction. In both cases, the realist inference moves from alternative counterfactual measurement possibilities to a single ontological description. In EPR, the move is synchronic: from the fact that the position of $A$ could be inferred if the position of $B$ were measured, and that the momentum of $A$ could be inferred if the momentum of $B$ were measured, Einstein and Popper conclude that $A$ possesses both position and momentum at one and the same time. Bohr denies that the two counterfactual attributions can be combined, since they belong to mutually exclusive experimental contexts. In Popper's analysis of measurement, the move is diachronic: from the fact that different measurements would allow different retrospective reconstructions of the particle's past, Popper infers a determinate trajectory through time. Hermann denies that these alternative reconstructions can be fused into one context-independent history. Thus Bohr resists a simultaneous state-attribution assembled from incompatible contexts, while Hermann resists a trajectory-attribution assembled from incompatible retrospective reconstructions.

%The disagreement can thus be located in the status of the conditional state-ascriptions licensed by alternative measurements. After the interaction, the two systems are no longer described by autonomous pure states, but by a joint entangled state. A measurement of the position of $B$ permits one to assign a corresponding position state to $A$; a measurement of the momentum of $B$ permits one to assign a corresponding momentum state to $A$. Einstein and Popper interpret these alternative conditional assignments as revealing properties that $A$ already possesses independently of the measurement performed on $B$. Bohr, and in an analogous way Hermann, deny this step. The two assignments belong to different experimental contexts and cannot be combined into a single context-independent description of $A$. Thus the disagreement does not concern merely whether both measurements can be performed in succession, but whether two counterfactual measurement contexts may be fused into one state attribution.

%In Popper’s reconstruction, EPR change the theory by treating the entangled systems as if they already had autonomous individual states before measurement. That is precisely what the quantum formalism does not give them.


%Popper accuses Bohr of changing the theory by refusing to combine the two EPR inferences. But one may reverse the charge. In Popper's reconstruction, it is EPR who tacitly go beyond the theory, since they treat the two entangled systems as if, before measurement, each already possessed an autonomous state of its own. Quantum mechanics does not license this assumption. After the interaction, the theory assigns a state to the compound system $A+B$, not separate pure states to $A$ and $B$. A measurement of the position of $B$ permits a conditional position-state to be assigned to $A$; a measurement of the momentum of $B$ permits a conditional momentum-state to be assigned to $A$. But these conditional assignments do not imply that $A$ possessed, prior to measurement and independently of context, a complete state containing both quantities. Thus the contested step is not Bohr's restriction on the theory, but EPR's attempt to extract from the entangled state a context-independent state of the subsystem.

%Although he does not see to have now, he is interested that Popper, we would call rejected Hermann's argument just like he rejected Bohr's reply to EPR. For EPR argument, we then obtain, for that instant of time, positions and momenta of both $A$ and $B$. Admittedly, the momentum of $A$ and the position of $B$ will have been destroyed or smeared by these measurements at a certain time: 

%What 
%
%
%
%
% if I can successfully carry out my decision without interfering with you—whether I wish to measure your size o r your weight, then it is unreasonable to say that you do not have both size and weight a t the same time but that you adopt the one or other b y adjusting yourself t o the stimulus of the measurement.

%\q{surreptitious change of the theory which Einstein attacked}.
%
%\q{This leaves completely open what the particle itself does}







%However, these retrodictions are not trajectory reconstructions, but as Hermann pointed out, measurement reconstructions.



%  \q{it is not possible, according to quantum theory, to project the path further back, \ie, to the region of time \myemph{before} the first of these measurements.}


%\q{non-predictive measurements determine the path of a particle only \myemph{between} two measurements}. 

%In particular, he recognized his intra-contextual mistake of his \TE: \q{it is not possible, according to quantum theory, to project the path further back, i.e., to the region of time \emph{before} the first of these measurements}. \autocite[242\fn{3*}]{Popper1959} . However, Popper did not seems to have to have appreciated deeper inter-contextual critique:  \q{non-predictive measurements determine the path of a particle only \emph{between} two measurements,} such as a measurement of momentum followed by one of position (or vice versa). 

%Without this, the possibility of reconstructing the path before the first momentum measurement Popper's \TE fails



%That the past tajectroies are not verifiabe, but they wihout them \qm could be falsifiable. \q{Measurements are retrodictions,} he wrote \citep[25]{Popper1967}. 

%Indeed, in The dogma of Copenhave interpation is that, ) those vectors in Hilbert space) insisted that an electron just cannot have a sharp position and momentum at the same time. However, that 

%Indeed, this notion will contiune to play a central role in his critique of the Copenehagen interpatio. Indeed, thesis $7$ of his interration of \qm got back to Heisenerg's \emph{Geschmackfrage} remark.   For Popper, this meant that context-free \emph{trajectory reconstructions}; the possibility of a  state is not a \emph{matter of taste}, but a \emph{matter of necessity}; they are needed to test the theory. Hermann's objections went completely unnoticed. Indeed, measurements are retrodictions, in the sense that quantum mechanics needs to \emph{measurement reconstruction}, and thse are a matter of context.





%Weizsäcker im Kontext

%Indeed, the debate did not change in later years. Both Popper and Hermann might have accepted Popper's phrase \s{measurements are retrodictions}. For Popper, such measurements are context-independent \emph{trajectory reconstructions}: they retrospectively determine the position-cum-momentum path of an individual particle and thereby test the statistical scatter relations. For Hermann, by contrast, retrodiction can only mean a context-dependent \emph{measurement reconstruction}: one reconstructs the outcome of a definite experimental arrangement, either for measuring position or for measuring momentum.


%The Opper: the was trajectro. howver, that  For Popper, the statistical content of quantum mechanics concerns the scatter of possible paths, while precise retrodictive measurements, bt can reconstruct individual trajectories and thus test the scatter relations. For Hermann, the statistical content of quantum mechanics concerns. \q{measurements are retrodiction}: For Popper, a retrodictive measurement reconstructs a context-independent \emph{trajectory reconstructins}: the path of For Hermann, by contrast, contexst depend \emph{measurement reconstructions}: from the spot on scren to  

%within a determinate experimental context: what is reconstructed is not an independently defined trajectory, but the outcome made meaningful by a particular measurement arrangement.
 

%For Popper, the statistical content of quantum mechanics concerns the scatter of possible paths, while precise retrodictive measurements can reconstruct individual trajectories and thus test the scatter relations. 
%
%For Hermann, the statistical content of quantum mechanics concerns the distribution of possible measurement outcomes within a determinate experimental arrangement; it does not license the reconstruction of a context-independent path specified by both position and momentum.



%the distribution of possible measurement outcomes within a determinate experimental arrangement; it does not license the reconstruction of a context-independent path specified by both position and momentum.


%Popper \q{these measurements are retrodiction}, btu trajeor recontru, that Hermann these measurements are retrodiction in the sese measuremetn reconstricitso.



%Einstein treats state-identity as prior to state-ascription: the second system has the same real state if it is not disturbed, and the question is whether the quantum state gives a complete description of that state. Bohr, by contrast, denies that state-identity can be fixed independently of the experimental arrangement. For him, the very conditions under which one may speak of the same physical state are supplied by the measurement context. Thus the issue is not merely that different arrangements license different state-ascriptions, but that there is no context-independent criterion of state-identity to which those ascriptions could be compared.

%Popper later conceded that intra-contextual mistake, that one cannot reconstruct the trajectory before measurement was conviced; howver, Popper never properply appreacitaed the deeper sense of Hermann's critique that very concept of past trajectroy is meanigless in \qm.



%But in Einstein’s generic argument, the contextual state-ascription may be only a preparation of some conditional state, not a measurement of a definite property of the second system. Thus the usual Bohr machinery loses part of its grip.

%\cop{Bohr can still say that the choice of measurement context changes the conditions for assigning a quantum state to the second system. But he cannot always supplement this with a concrete reconstruction of a measured eigenvalue}.

%In the EPR case, Bohr can give a positive account of the experimental situation: the measurement record, interpreted within the whole arrangement, allows the experimenter to reconstruct a definite aspect of the interaction and hence to infer the value of a determinate observable. In the generic entangled-state case, however, this positive account breaks down. The conditional states \( |\psi_n\rangle \) and \( |\varphi_s\rangle \) assigned to the second system need not be eigenstates of any observable of that system. Thus there need be no eigenvalue to which the record can be traced back. Bohr may still insist that the meaning of state-ascription is context-dependent, but he no longer has the same explanatory resource that he has in the EPR case: a reconstruction from the experimental record to a definite measured quantity.

%That is why Bohr’s positive account is less effective against Einstein’s generalized version. Bohr can still say that the choice of measurement context changes the conditions for assigning a quantum state to the second system. But he cannot always supplement this with a concrete reconstruction of a measured eigenvalue or a definite property of the second system. In the general case, the “non-mechanical disturbance” affects the state assignment itself, not necessarily the possibility of reconstructing a determinate value of some observable.

%With the exception of Schrödinger, early respondents to the EPR paper were not aware of the difference between Einstein’s private argument and the published EPR argument.  %For Bohr, complementary measurements do not yield two incomplete descriptions of one and the same phenomenon. Since the experimental arrangement is an integral part of the phenomenon, mutually exclusive position and momentum arrangements define different phenomena.

% For Bohr, complementary measurements do not yield two incomplete descriptions of one and the same phenomenon. Since the experimental arrangement is an integral part of the phenomenon, mutually exclusive position and momentum arrangements define different phenomena. Indeed, it was Pauli that seems precisely sensed this point.

%Howver, Bohr something stronger, or at least something you can reconstruct as stronger. A decade later, I read Einstein's argument in the latter's short \citeyear{Einstein1948} paper for a special issue of \emph{Dialectica} that Pauli had curated. In the editorial, \citet{Pauli1948} effectively shifted the notion of quantum \s{state}. Pauli pushed the contextual argument a step forward. For Einstein, a physical real state must not depend on the choice of the experimental arrangement. For Pauli, by contrast, the quantum-mechanical state is defined within a context that includes the experimental arrangement. Not only is the state attributed to $B$ dependent on the experimental context; the very identification of what counts as the \s{same state} across different contexts is not available independently of that context. If the composite system $AB$ is described by a pure entangled state, then neither subsystem need possess its own autonomous pure $\psi$-function. This means, however, that before measurement neither $A$ nor $B$ had its own autonomous pure state in Einstein’s sense. the electron cannot be assumed to remain in \s{the same state} independently of how the photon is measured, because the electron has no real state to begin with.


%which was exactly what Einstein found unbearable: 

%then the complete description of the real state of the system,
%whatever it might be, must be independent of the choice of measurement of A or B

%two pictures are complementary to each other, in the sense that either picture is separately applicable, and provides a physical story that explains what happens in the measurements on the two sides, but the two pictures cannot be combined into one.



%physical reality context, physical reality 
%This is nothing but the context-dependence that is indeed \citep{Henry-Hermann1948}, one of the Copenhagen interpretations \citep{Pauli1950}. Whereas Popper will be come in the, Hermann remained one of this will be remembered.



%It was this  that Popper later mistook as subjecti

%This dependence does not mean, however, that the state remains arbitrarily subject to the experimenter’s influence. Rather, once the experimental arrangement has been \emph{chosen}, the physical state can be regarded as partially detached from that influence. 

%



%Sie gilt vielmehr nur relativ zu den Beobachtungen, auf Grund deren sie aufgestellt
%worden ist, und läßt die Möglichkeit offen, das gleiche System, und zwar für den
%gleichen Zeitpunkt, relativ zu einem andern Beobachtungszusammenhang durch
%eine andere Wellenfunktion zu beschreiben. Im Beispiel der skizzierten Orts-
%messung: Relativ zum alten Beobachtungszusammenhang, in dem Elektron und
%Lichtquant je durch die eigene Wellenfunktion gegeben waren, führt die Wech-
%selwirkung zwischen beiden zu jener Wellenfunktion, die sich auf beide Elemen-

%The issue at stake in the Popper--Hermann debate was the possibility of context-free trajectories, discussed under the heading of causality. Yet this issue already pointed toward a more general problem: whether a physical system can be assigned a context-free state at all.\todo{??} 




%If some of the quantities that can be predicted for $B$ are mutually incompatible, such as position and momentum, then, Einstein argues, these quantities must nevertheless both belong to the real state of $B$, since either prediction can be obtained by a free choice of measurement on $A$ without disturbing $B$. 

% Thus the issue is not only that different $\psi_B$-functions can correspond to the same state, but also that the same state must contain more determinate physical content than any one of those wave functions represents.

%Pauli can therefore concede that, after the result on $A$ has been registered, $B$ is no longer assigned the \emph{same} conditional state. Yet this does not imply any physical action at a distance, since the marginal statistical predictions for observations on $B$ considered in isolation has remained the same, and only That issue at the Popper-Hermann debate on possiblit context-free trajectory was a to more general question context-free physical state. 


%^ but a change in the state attribution defined by the phenomenon as a whole\footnote{What remains unaffected are  Thus Pauli can reject Einstein’s demand for a context-independent real state of $B$ without abandoning the requirement that no controllable influence is transmitted from $A$ to $B$}. 


%Hermann and Popper had discussed the problem of trajectories in the context of causality. Popper’s context-free trajectory, which Hermann rejected as impossible, The later EPR discussion shows, however, that the underlying issue was more general: the problem of trajectories is a special case of the problem of state attribution.  Einstein’s context-free state express the same classical demand: the physical description of a system should be identifiable independently of the experimental arrangement through which it is reconstructed. The Bohr--Pauli reply denies this demand. 

%Both trajectory attribution and state attribution are meaningful only within a phenomenon, that is, within the experimental conditions that define the possible description. %Popper’s trajectory is a context-free dynamical state: he assumes that one can combine determinations belonging to different experimental contexts into a single history of the particle. Einstein’s separability argument generalizes the same demand: the distant system $B$ must possess a context-free physical state, whose identity is independent of the experimental arrangement chosen for $A$.


%It requires only a \emph{criterion of sameness}. If system $B$ is no longer interacting with system $A$, then the physical state of $B$ should remain the \emph{same}, whatever measurement is freely chosen on $A$. But quantum mechanics assigns different $\psi_B$-functions to $B$ depending on that choice. The difficulty is therefore not, in the first instance, that $B$ must possess both position and momentum. It is rather that the same physical state of $B$ appears to receive different wave-function descriptions. At this point, the Hermann--Bohr reply must become more radical. It is not enough merely to reject the inference from prediction to a determinate real state of $B$. One must ask whether it is legitimate to speak, across the two experimental contexts, of the same real state of $B$ in the sense presupposed by Einstein.



%Pauli’s ultimate reply is to deny Einstein’s assumption that there is a context-independent sense in which $B$ has \s{the \emph{same} state} across the different possible measurements on $A$.


%The concept of state, fundamental in quantum mechanics, is defined by a complete account of the reproducible experimental conditions to which the observed system is subjected. Pauli identified the \s{state} of a system with Bohr’s notion of \s{phenomenon}: an correspondence between an indicator reading, an experimental apparatus, and an eigenvalue. 


%toward Bohr’s notion of \s{phenomenon}. Bohr does not seem to have been fully aware of the difference between Einstein’s private argument and the EPR argument. It was Pauli, however, who, a decade later, responded along these lines once he had read the argument in Einstein’s short \citeyear{Einstein1948} paper for a special issue of \emph{Dialectica} that Pauli had curated. In the editorial, \citet{Pauli1948} effectively shifted the notion of quantum \s{state} 


%Pauli appealed to Bohr’s notion of \s{phenomenon}: it is the phenomenon, in this sense, that defines the state of the observed system. The concept of state, fundamental in quantum mechanics, is defined by a complete account of the reproducible experimental conditions to which the observed system is subjected. In this, sense are indeed, in differnt state; 


%Pauli’s ultimate reply is to deny Einstein’s assumption that there is a context-independent sense in which $B$ has \s{the \emph{same} state} across the different possible measurements on $A$. The disagreement turns on the very definition of \s{state}. As Pauli explicitly recognized, Einstein’s position presupposes the existence of an objective \s{real} state of a physical system that is independent of any selections and results of experimental arrangements that do not actually disturb the system under consideration. Against this, Pauli appealed to Bohr’s notion of \s{phenomenon}: it is the phenomenon, in this sense, that defines the state of the observed system. The concept of state, fundamental in quantum mechanics, is defined by a complete account of the reproducible experimental conditions to which the observed system is subjected. Thus the state attributed to $B$ is not an autonomous property of $B$ considered apart from the experimental context. It is fixed only within a phenomenon, that is, within a context that includes the chosen experimental arrangement and the conditions under which the relevant prediction or attribution is made.





%For Bohr, complementary measurements do not yield two incomplete descriptions of one and the same phenomenon. Since the experimental arrangement is an integral part of the phenomenon, mutually exclusive position and momentum arrangements define different phenomena. Bohr does not seem to have been fully aware of the difference between Einstein’s private argument and the EPR argument. It was Pauli, however, who, a decade later, responded along these lines once he had read the argument in Einstein’s short \citeyear{Einstein1948} paper for a special issue of \emph{Dialectica} that Pauli had curated. In the editorial, \citet{Pauli1948} effectively shifted the notion of quantum \s{state} toward Bohr’s notion of \s{phenomenon}. A description of the state is complete when it establishes a unique reconstructive chain linking the macroscopic registered outcome on some indicator, the experimental arrangement, and the corresponding microscopic eigenvalue.

%Pauli’s ultimate reply is to deny Einstein’s assumption that there is a context-independent sense in which $B$ has \s{the \emph{same} state} across the different possible measurements on $A$. This was ultimately the point of the discussion. On Einstein’s view, the attribution of a physical state must be independent of the experimental arrangement: an objectively existing physical state should be fully detached from the experimenter’s choice of how to investigate it. Quantum theory, by contrast, requires a lower degree of detachment. The state attributed to a system depends on the experimenter’s free choice between mutually exclusive experimental arrangements. This dependence, however, does not make the state attribution arbitrary or merely subjective. Once the experimental arrangement has been chosen, the attribution is fixed by objective constraints: the arrangement itself, the registered outcome, and the quantum-mechanical rules that connect them.




%Pauli’s reply, equaite the notion of \s{state} with the notion of \s{phenomenon}, appears to deny even this weaker assumption. The different possible measurements on $A$ do not merely provide different descriptions of one and the same independently identifiable state of $B$. They belong to different phenomena, thus they are indeed different states, each defined by its own experimental arrangement. The disagreement therefore turns on the definition of \s{state}. For Einstein, a physical state must be independent of the experimental arrangement: it is the state of an objectively existing system, detached from the experimenter’s choice of how to investigate it. For Bohr and Pauli, by contrast, the quantum-mechanical state is defined only within a phenomenon, that is, within a context that includes the experimental arrangement by which the relevant prediction or attribution is made.

%Pauli can then concede that an observation on $A$ changes the state attributed to $B$, while denying any physical action at a distance, thereby preserving the notion of a \textit{closed system}. After the interaction, one can decide to measure, say, position or momentum on $A$. The marginal statistics for $B$, considered in isolation, remain unchanged; what changes are the conditional probabilities determined by the observed result on $A$. Before a selective observation on $A$, the theory assigns probabilities to possible outcomes for $B$; after $A$ is found, for example, at position $q$, the conditional probability for the correlated position of $B$ may become unity. The change is therefore not a local disturbance of $B$, but a change in the context-relative state assignment. 



%



%This dependence does not mean, however, that the state remains arbitrarily subject to the experimenter’s influence. Once the experimental arrangement has been chosen, the corresponding state attribution is fixed by objective constraints: the arrangement, the registered outcome, and the quantum-mechanical rules linking them. The state is therefore not detached from the experimental context as Einstein required, but it is detached from the experimenter’s further arbitrary intervention within that context.

\todo{Tis wheter depnde on the xperimeatl arramgent or whter id ote not.}

%The historical line reconstructed in this paper thus leads from Popper’s failed trajectory reconstruction to Einstein’s separability argument, and from Hermann’s measurement reconstruction to the later Bohr–Pauli response. Popper tried to use retrodiction to recover an underlying classical history behind the statistical formalism. Einstein shifted the issue to a deeper question: whether the same distant physical state can be represented by different wave functions. Hermann’s analysis did not yet answer this question in its later EPR form, but it supplied an important part of the conceptual machinery needed to do so. It showed that state attribution in quantum mechanics is not a context-independent act of naming an underlying property; it is a reconstruction tied to an experimental arrangement. The same lesson, reformulated through Bohr’s notion of phenomenon and sharpened by Pauli, became one of the central replies to Einstein’s separability argument.

%Einstein’s incompleteness argument then reappears as the claim that a theory which only secures statistical closure, without describing the individual separated system, is not a complete theory of physical reality.

%

%just before \citets{Schroedinger1935} paper 



%
%
%
%
%%The public replies to EPR, such as \citets{Bohr1935}, typically challenged the use of the \emph{criterion of reality}: from the possibility of predicting a quantity with certainty, without mechanically disturbing the system, it does not follow that this quantity belongs to the system as an independently real property. On Bohr’s reading, the very conditions under which such a state attribution is possible depend on the experimental arrangement.
%
%
%
%
%
%%and th We shall use Hermann’s analysis to shine further light on both Heisenberg’s and Bohr’s replies to EPR. 
%
%%her analysis can be read as pointing in this direction. It is perhaps not accidental that \citet{Heisenberg1935} used \vW's $\gamma$-ray microscope in his unpublished response to EPR \citep[**]{BacciagaluppiCrull2024}.
%
%
%%quantum mechanical state description of a physical system 
%%
%%state description in classical physics i
%
%%They belong to mutually exclusive experimental arrangements. If one measures the position of the first system $A$, one can infer the position of $B$; if one measures the momentum of $A$, one can infer the momentum of $B$. Yet these are counterfactual alternatives, not jointly realizable determinations.
%
%%Hermann expresses there once again with great clarity the decisive lesson the fundamental issue at stake: quantum-mechanical \emph{state attribution}, including that  \q{is always relative to the context of observation} \citep[**]{Hermann1936}. The prediction of position and the prediction of momentum  for the system $B$ are not two determinations made within one and the same experimental context. However, as we now know \citep{Howard1985}, Einstein’s private version of the argument was designed to carefully avoid state attribution of the distant system $B$. It requires only a  \emph{state identity} is independent of the context of observation. Against this argument, the Bohr-Hermann response is less pressing: there is not definite property of the second system, thus not measurement, and no measurement context\footnote{However, Hermann could possibly have thought that even state identity is context-dependent. In Hermann's \TE, the electron cannot be assumed to remain in \s{the same state} independently of how the photon is measured, because the electron has no real state of its own prior to measurement. The consequence was possibly more radical than Hermann envisaged: The electron--photon pair is not treated as two distinct objects, each with its own separate state, but rather as a single object with one indivisible state}. 
%
%
%%Popper acknowledges that the conclusion is inescapable: \q{non-predictive measurements determine the path of a particle only \myemph{between} two measurements,} such as a measurement of momentum followed by one of position (or vice versa): \q{it is not possible, according to quantum theory, to project the path further back, \ie, to the region of time \myemph{before} the first of these measurements.}\autocite[242\fnstar{3}\me]{Popper1959} quantum-mechanical \emph{state attribution}, including that  \q{is always relative to the context of observation} \citep[**]{Hermann1936}.


%Thus the focal plane is “out of focus” with respect to the object, but “in focus” with respect to direction. That is exactly why it yields momentum information rather than position information.


%Taking O as the x-axis and the optical axis as the y-axis, the scattered photon travels mostly upward, so its momentum has a large y-component and, if it is slightly inclined, a smaller x-component:

%The focal-plane registration determines the direction of the photon, hence $\vartheta_{,}$and therefore its transverse component $p_{\gamma, x}$. It is this transverse component, not the upward component $p_{\gamma, y}$, that is relevant for determining the electron's momentum along $L$.

%In the focal-plane arrangement, by contrast, the lens maps photons with the same outgoing direction onto the same point of the focal plane. Since the outgoing direction determines the transverse momentum component of the photon, the registration on $S$ is interpreted as momentum information and, by conservation of momentum, as information about the electron's $p_x$. It does not determine the electron's position $x$.

%In the image-plane arrangement, the scattering point $P$ functions as the source of a spherical wave. The lens collects any  from different directions within its aperture and recombines them at the image point $P'$ on $S$. The larger the aperture, the larger the portion of the spherical wave that enters the lens, and the more sharply the apparatus can distinguish a wave originating at $P$ from one originating at a neighboring point. A wider aperture therefore yields a more precise determination of the electron's position $x$.

%focal plane any photon reaching the lens with the same outgoing direction is focused onto the same point of the focal plane.

%within the cone accepted by the lens.
%A photon is scattered by the electron at some point $P$ in the object plane $L$. After the scattering event, the photon travels upward into the microscope, 


% In Hermann's terms, they license either a position or a momentum \emph{measurement reconstruction}, but not a \emph{trajectory reconstruction}. Hermann therefore concludes that Heisenberg's matter-of-taste remark was unfortunate: \qt{Because of this contrast between quantum-mechanically admissible retrodictions and fictitious trajectory calculations, I think that Heisenberg's much-disputed statement about these trajectory calculations, namely that it is a mere matter of taste whether one should attribute physical reality to them, must be corrected as follows: these calculations leave out of account the relation of every physical determination to the relevant observation context, which is decisive in quantum mechanics, and in this respect become physically meaningless}{Wegen dieses Gegensatzes zwischen quantenmechanisch zulässigen Rückschlüssen und der fiktiven Bahnberechnungen meine ich, dass man die viel umstrittene Äusserung Heisenbergs über diese Bahnberechnungen: es sei eine reine Geschmacksfrage, ob man ihnen physikalische Realität zuordnen solle, dahin korrigieren muss, dass diese Rechnungen die quantenmechanisch entscheidende Beziehung jeder physikalischen Bestimmung auf den jeweiligen Beobachtungszusammenhang ausser Acht lassen und insofern physikalisch bedeutungslos werden} \letterGHp{Hermann}{Weizsäcker}{5}{3}{1935}[14].


%As a curiosity, and to characterize Popper, here is a line of thought which, according to Popper, underlies the uncertainty relations: the energy exchange occurring in a measurement between object and apparatus changes the state of the object, so that after the measurement it is different from before. Thus, through measurement, one always learns a state that has just been destroyed by the measurement. Therefore one cannot infer the state of an atomic object after the measurement from the measurement itself, and the measurement cannot be used as a basis for predictions. This is intelligible only if one assumes that Popper regards precisely the fictitious trajectory calculations as the only complete measurements.

%\vW adds that Popper’s book had already been printed before Popper saw his reply, and that Popper’s subsequent replies had so far been answered privately by letter. He nevertheless regards Popper as reasonable insofar as he gives up views once they have been clearly refuted.


% His electron-diffraction example makes this point explicit: if an electron is later found in one of the directions favored by the wave picture, one may say that such a deflection would have been too improbable without the electron’s wave properties. But this already introduces probability, and probability becomes meaningful only through the performance of many experiments. Similarly, the idea of proving the wave character of light by



%As we shall see, that will go to the very core of the question, beund \vW. 

%He cannot think of a single, individual experiment that would count as proof of the wave picture. 





%For \vW, moreover, an “impulse selection” is not different in principle from an “impulse measurement,” because of the symmetry of nature with respect to past and future.


%Howver,  \q{Ob diese „Anweisung, Herrn Popper zu widerlegen“ in einen zwingenden allgemeinen Gedankengang umgeformt werden kann, weiss ich nicht recht. Letzten Endes wird man sich eben doch auf den Formalismus im ganzen berufen müssen, den Herr Popper nicht verstanden zu haben scheint.}




%All particles arriving at a point $P$ would then have the same momentum.  \VW objects that $y$ cannot be known with arbitrary precision. First, the beam cannot be made narrower than the aperture width $d$; second, if $d$ is made too small, diffraction occurs, and the wave picture must be used. The attempt to sharpen the position selection therefore introduces an indeterminacy in the particle's momentum, and hence in its velocity during the measurement process. In this version of the argument, the mutual exclusiveness of the position and momentum measuring precudes is more explicit.


%\VW's point is that Popper's selection of electron momenta in an electric field has the same structure as the field-deflection measurement discussed by Heisenberg for a magnetic field. The deflection angle can indeed be used to infer the particle's velocity or momentum. But the angular determination presupposes a beam defined by apertures, and these apertures introduce precisely the diffraction and localization uncertainties that restore the uncertainty relation. Popper therefore treats the field-deflection apparatus as if it yielded a sharp momentum selection without paying the corresponding price in spatial indeterminacy.

%The consequence is that the arrangement cannot do what Popper wants. It cannot support a precise inference from a known position before the measurement to a position after the measurement, as in Heisenberg’s discussion of magnetic-field momentum measurement; nor can it support Popper’s reverse inference from the later position, namely the time of arrival at $P$, back to the earlier position. In both directions, the attempt to combine precise momentum selection with precise positional information is defeated by the finite duration of the measurement and by diffraction effects at the aperture.

%One possible realization would be resonance fluorescence: a gas transmits all light except for one spectral line, which, as Wood had shown experimentally, is reflected almost like light from a metallic surface. In this sense the gas can be said to filter out light quanta with one definite momentum. Thus Popper's argument fails for the same reason as the procedure analyzed by Heisenberg: the attempt to combine precise trajectory reconstruction with precise momentum selection is blocked by the uncertainty relations.


%\vW's point is that Popper's selection of electron momenta in an electric field has the same structure as the field-deflection measurement discussed by Heisenberg for a magnetic field. The deflection angle can indeed be used to infer the particle's velocity or momentum. But the angular determination presupposes a beam defined by apertures, and these apertures introduce precisely the diffraction and localization uncertainties that restore the uncertainty relation. Popper therefore treats the field-deflection apparatus as if it yielded a sharp momentum selection without paying the corresponding price in spatial indeterminacy.

%\VW then applies the same point to Popper's proposed momentum selection of electrons in an electric field, summarizing the objection he had already made in private correspondence with Popper. \VW's point is that Popper's selection of electron momenta in an electric field has the same structure as the field-deflection measurement discussed by Heisenberg for a magnetic field. Popper assumes that particles with different momenta are deflected by different amounts, so that the displacement $y$ of a deflected particle from an undeflected one directly indicates its momentum. 


%In both cases, the central idea is that particles with different momenta are deflected by different amounts, so that the observed displacement appears to provide a direct measure of their momentum. Popper therefore assumes that all particles arriving at the point $P$ have the same momentum.


%

%Every bundle of rays that entered the lens parallel to each other is perfectly focused into a single, sharp point on the focal plane.


% At the image plane $d_i$, the rays that originated from the same point in the object plane have been brought together; at the focal plane $f$, the rays that entered the lens with the same direction have been brought together. 

% By choosing where to place $S$, one selects which sorting criterion to intercept and record. 


%By definition, a convex lens is designed to take any bundle of light rays traveling parallel to its optical axis and bend them so they all intersect exactly at the focal point, located at the distance $f$. 

%Now, consider an object placed at a finite distance $d_o$ in front of the lens. Because the object is close, the light rays originating from any single point on it are actively spreading apart (diverging) when they hit the glass.

%This classical optical alternative constrains the operation of the \HW microscope, which functions as a selection or preparation device for incoming photons either according to position or according to direction. Now, consider an object placed at a finite distance $d_o$ in front of the lens. Because the object is close, the light rays originating from any single point on it are actively spreading apart (diverging) when they hit the glass.

%When light originates from an object that is infinitely far away, the rays have spread out so much by the time they reach the lens that they are traveling perfectly parallel to one another. They possess zero divergence.

%Because the curvature of the glass and its refractive index do not change, the lens can only bend a light ray by a specific, fixed angle.  It applies this exact same inward "push" to whatever light strikes it, regardless of where that light came from


%By definition, a convex lens is designed to take any bundle of light rays traveling parallel to its optical axis and bend them so they all intersect exactly at the focal point, located at the distance $f$.

%same angle $\alpha$


%The registration therefore selects a momentum context.


%If a bundle of rays enters the lens parallel to the optical axis, it has no initial transverse divergence. The lens-induced deflections are then just sufficient to make the rays intersect at the focal distance $f$

%One may think of the lens, by analogy, as giving each ray an inward transverse deflection whose magnitude depends on where the ray crosses the lens: rays farther from the optical axis are bent more strongly than rays closer to it. If a bundle of rays enters the lens parallel to the optical axis, it has no initial transverse divergence. The lens-induced deflections are then just sufficient to make the rays intersect at the focal distance $f$.
%
%If, by contrast, a bundle of rays reaches the lens after diverging from a common point $P$ in the object plane $O$, the rays already possess an outward angular spread. The lens still bends them inward, but part of its action first compensates for this prior divergence. The resulting convergence is therefore weaker, and the rays meet only farther away, at the image distance $d_i$. Thus the focal plane is where parallel incoming bundles are brought to coincidence, whereas the image plane is where bundles diverging from a common object point are brought to coincidence.

%One may think of the lens, by analogy, as giving each ray an inward transverse deflection whose magnitude depends on where the ray crosses the lens: rays farther from the optical axis are bent more strongly than rays closer to it. If a bundle of rays enters the lens parallel to the optical axis, it has no initial transverse divergence. The lens-induced deflections are then just sufficient to make the rays intersect at the focal distance $f$.
%
%If, by contrast, a bundle of rays reaches the lens after diverging from a common point $P$ in the object plane $O$, the rays already possess an outward angular spread. The lens still bends them inward, but part of its action first compensates for this prior divergence. The resulting convergence is therefore weaker, and the rays meet only farther away, at the image distance $d_i$. Thus the focal plane is where parallel incoming bundles are brought to coincidence, whereas the image plane is where bundles diverging from a common object point are brought to coincidence. %By contrast, $\theta$ denotes the angle between the refracted ray after the lens and the vertical axis. This latter angle characterizes the geometrical focusing of the ray and should not be identified with the photon's pre-lens propagation direction.

%In this case, however, the point in the object plane from which the light originated, and hence the position at which the collision occurred, remains indeterminate. The same physical interaction must therefore be characterized differently: the electron is now described by approximatly a plane wave with relatively sharp momentum but indeterminate position. Narrowing the effective angular acceptance improves the selection of the photon direction, and hence of the momentum transfer, but at the price of losing information about the point of origin.




%This asymmetry in reception reveals the complex and evolving relation between philosophy and the sciences, and physics in particular. Popper's technically fragile and ultimately misguided reflections on \qm were embedded in a broader philosophical vision that transformed certain Kantian-Friesian motifs into a historically consequential account of scientific rationality. Hermann's Kantian-Friesian framework was just as ambitious as Popper's, but it did not survive as a compelling general philosophy of science, carrying with it her penetrating and technically informed interpretation of \qm. Popper's philosophy became enormously influential, while Hermann's analysis was largely forgotten. However, the tables seem to have turned. Grand synthetic philosophies of science have come to appear suspect, and Popper's philosophy, among others, has fallen out of fashion. At the same time, Hermann's analysis has acquired new visibility, as philosophy of physics has become increasingly technical. Yet the tables may turn again, as philosophy of the special sciences seems increasingly unable to connect local technical questions to philosophical issues of wider significance.

%This suggests a possible mediating role for the history of science. For a single philosopher, it is increasingly difficult to remain technically up to date with contemporary research programs such as quantum gravity while also connecting them to broader questions in ethics, politics, or aesthetics. Historical theories such as relativity and \qm offer a different kind of opportunity. They are old enough for their mathematical, experimental, and conceptual structures to be studied with sufficient depth, yet still philosophically alive enough to raise questions about causality, objectivity, realism, determinism, and the nature of scientific rationality. The history of science can therefore mediate between local technical analysis and global philosophical reflection. It allows one to recover the philosophical stakes of scientific theories without sacrificing attention to their formal and experimental details.

%philosophy of physics risks losing even the ambition to articulate a philosophical vision of physics as a whole, not to mention of science.

% narrowness of much professional Anglophon Popper had the global philosophical ambition but lacked sufficient local technical control in \qm. Hermann had the local technical penetration, but her global framework did not remain persuasive. 

%Popper had the global philosophical ambition but lacked sufficient local technical control in \qm. Hermann had the local technical penetration, but her global framework did not remain persuasive.

%Die prinzipiellen Schranken der Vorausberech-
%nung werden dadurch und nur dadurch verständlich, daß die quantenmechanische
%Naturbeschreibung sich als relativ erweist, als relativ zum jeweiligen Beobach-
%tungszusammenhang, in dem der Forscher sich seinem Objekt gegenüber befindet.
%Dieser relative Zug der Naturbeschreibung ist der klassischen Physik fremd; für
%sie ist die Charakterisierung irgendeines Systems eindeutig und unabhängig von
%der Art, wie der Beobachter sich Kenntnis von ihm verscha&. Und daher kommt
%sie zwingend zu der Annahme, daß bei hinreichender Beobachtungsschärfe und
%hinreichender Kenntnis der naturgesetzlichen Zusammenhänge die Untersuchung
%physikalischer Systeme die Ursachen ihrer weiteren //125// Entwicklung mit belie-
%biger Deutlichkeit zu bestimmen und damit diese weitere Entwicklung vorauszu-
%berechnen gestattet.


%Ultimately, the reconstruction of a trajectory would require a repetition of \WH setup for every instant $t$: at each $t$, one would have to reconstruct both the position $\pos(t)$ and the momentum $\mom(t)$ of the same individual system. But Hermann's argument shows that each such reconstruction is context-dependent. A given experimental arrangement may license the attribution of \pos or of \mom at time $t$, but not both together as elements of a single context-independent state description. Thus the retrospective reconstruction of one measured value at one time cannot be generalized into the reconstruction of a classical trajectory. 

%The trajectory calculation ignores the quantum-mechanically decisive fact that every physical determination is relative to a specific observational context. Put in slightly modern terms, one may choose a $q(t)$-history or a $p(t)$-history, but not a hybrid history that treats both as simultaneously valid descriptions of the same individual process \citep{Griffiths2003}.





% The crucial point is that one observational context determines which features of the wave and particle pictures may be used together. This is why the uncertainty relations are preserved. A trajectory reconstructed by combining data from different contexts is therefore not a genuine physical description. To the extent that it exceeds the uncertainty limits, it is physically empty: it cannot be controlled experimentally and provides no basis for prediction.


%In the microscope example, the propagation of light through the microscope is described in wave terms, whereas the interaction between light and electron must be described corpuscularly, as a collision with a light quantum. How these two descriptions are reconciled depends on the experimental arrangement. Thus the single observational context established by the registration on the photographic plate determines which features of the two pictures may legitimately be used. In this way, the uncertainty relations are preserved




%If the light is registered in the image plane of the observed object, the wave description treats it as a spherical wave spreading out from a point; correspondingly, the corpuscular description assigns the collision a sharp position but an indeterminate momentum transfer. If, by contrast, the observation is made in the focal plane of the microscope, the relevant wave description is that of a parallel bundle of rays; correspondingly, the corpuscular description assigns the collision a determinate momentum transfer but an indeterminate position. 


%Die Bedeutung dieses Gegensatzes ergibt sich aus dem nur relativen Charakter der quantenmechanischen Beschreibungsweise. Da jede physikalische Beschreibung und Erklärung von Vorgängen nur relativ zum jeweiligen Beobachtungszusammenhang gilt, so bleibt die Berechnung etwa einer Korpuskelbahn, die die Bestimmungen //123// verschiedener Beobachtungszusammenhänge zu einer Darstellung zusammenfaßt, eben in dem Maß, wie sie über die Unbestimmt- heitsrelationen hinausgeht, physikalisch leer: Sie ist unkontrollierbar und gibt keine Anhaltspunkte für Vorausberechnungen.

%Diese nachträgliche Konstruktion einer Elektronenbahn unterscheidet sich
%von den Rückschlüssen, die zur Interpretation einer Messung erforderlich sind, dadurch, daß sie dem Dualismus von Wellen- und Partikelbild keine Rechnung trägt, sich auf die anschaulichen Merkmale eines dieser Bilder beschränkt und, da jede einzelne Beobachtung auf beide Bilder Bezug nimmt und daher //122// keines von ihnen unbeschränkt anzuwenden gestattet, diese Merkmale durch zwei verschiedene Beobachtungen festlegt.  Auf diesem Weg gelingt es anscheinend, die zu verschiedenen Beobachtungszusammenhängen gehörenden Bestimmungsstücke in einer einzigen Beschreibung des physikalischen Systems zu vereinigen.

%Im Gegensatz dazu muß bei der Interpretation einer Messung der Kausalschluß, der aus dem Zustand des Meßinstruments den des beobachteten Objekts entnimmt, den Dualismus berücksichtigen. Denn für den Vorgang der Messung ist es wesentlich, daß jeder atomare Vorgang auch Wellencharakter, jede Wellenbewegung auch korpuskularen Charakter zeigt. Die eine Beobachtung des Meßapparats, die dem Rückschluß auf das Meßobjekt zu Grunde liegt, bestimmt dabei, in welcher Weise Wellen- und Partikelbild einander bei der vorliegenden Interpretation begrenzen.


%\q{So wird in dem betrachteten Beispiel die Ausbreitung des Lichts im Mikroskop im Wellenbild beschrieben, während für die Wechselwirkung zwischen Licht und Elektron die korpuskulare Vorstellung des Lichtquants herangezogen werden muß. Wie beide Vorstellungen miteinander in Einklang gebracht werden, hängt von der Art der Messung ab: Wird das Licht in der Bildebene des beobachteten Objekts aufgefangen, so wird man im Wellenbild mit der Vorstellung einer von einem Punkt sich ausbreitenden Kugelwelle arbeiten und dementsprechend dem korpuskular gedeuteten Zusammenstoß zwischen Elektron und Lichtquant einen scharfen Ort aber einen verschmierten Impulsaustausch zuschreiben. Stellt man die Beobach- tung in der Brennebene des Mikroskops an, so hat man es mit einem parallelen Strahlenbündel und demnach im Korpuskelbild mit einem genau bestimmten Impulsaustausch aber einem unscharfen Ort zu tun. Der eine Beobachtungszusammenhang, in den der Physiker durch die Beobachtung der photographischen Platte eintritt, bestimmt also, welche Merkmale beider Bilder Verwendung finden. Die Wahrung der Unbestimmtheitsrelationen ist damit gesichert.





%Die umstrittene Darstellung, die HEISENBERG diesem Sachverhalt gegeben hat, es sei eine reine Geschmacksfrage, ob man einer solchen Bahnberechnung physikalische Realität zuordnen solle, ist demnach dahin zu berichtigen, daß diese Rechnung die quantenmechanisch entscheidende Beziehung jeder physikalischen Bestimmung auf den jeweiligen Beobachtungszusammenhang außer Acht läßt und insofern physikalisch bedeutungslos wird. Umgekehrt hängt die Möglichkeit einer, wenn auch nur mittelbaren Kontrolle der Kausalbehauptungen, die bei der Interpretation eines Messungsvorgangs auftreten, daran, daß diese Interpretation dem Dualismus und damit implizite dem relativen Charakter der Quantenmechanik gerecht wird.


%The measurement outcome is therefore not the mark alone, but the mark together with the optical history that connects the scattering event to its registration on the plate.


%Hermann’s point is therefore not merely that the act of mesurement disturbs the object. It is that a registration event has no determinate physical significance in isolation. A mark on the photographic plate is not, by itself, either a position measurement or a momentum measurement. Its significance depends on the optical arrangement that connects the mark with the preceding scattering event. If the plate is placed in the image plane, at the distance $d_i$ determined by the object plane and the thin-lens equation, the lens maps object points onto image points; the mark is then interpreted as position information. If the plate is placed in the focal plane, at the distance $f$, the lens maps directions of propagation onto points on the plate; the mark is then interpreted as momentum information. The measurement outcome is therefore not the mark alone, but the mark as embedded in a definite optical reconstruction.


%This also clarifies Popper's thought experiment. After the electron--photon collision, the two particles are described by the kind of joint state analyzed by Hermann. If one chooses to measure the momentum of the electron, the conservation law fixes the corresponding momentum of the photon. But this choice selects the momentum representation of the joint state. The electron's position is therefore not simultaneously fixed; it is given, at most, by the probability distribution associated with a possible position measurement. Popper's mistake is to treat the momentum inferred from the photon and the later position measurement on the electron as if they could be combined into a single trajectory. Hermann's point is that the first measurement determines only a momentum correlation within a momentum context. It does not also determine the electron's position history.


%The same joint wave function of the electron--photon system can be expanded in different systems of eigenfunctions. If one decides to set $S=f$, one expands it in eigenfunctions of the momentum operator, one obtains a representation in which each term correlates a photon momentum eigenfunction $\psi^P_p$ with the corresponding electron momentum eigenfunction $\psi^E_{p}$, according to the conservation law $p_E+p_P=P$. If decides to set $S=d_i$, one expands the same wave function in eigenfunctions of the position operator, one obtains instead a representation in which each term correlates a photon position eigenfunction $\psi^P_x$ with the corresponding electron position eigenfunction $\psi^E_x$, according to the relevant geometrical correlation, for instance $x_E-x_P=0$. Thus the same quantum-mechanical state may be represented either in terms of momentum correlations or in terms of position correlations; but these two expansions correspond to different experimental arrangements and cannot be combined into a single context-independent attribution of both position and momentum to the electron at the same time $t$.



%If $M$ is the magnification of the microscope, this correlation may be written schematically as
%\begin{equation*}
%x_S=Mx_E.
%\end{equation*}
%Thus the same quantum-mechanical state may be represented either in terms of momentum correlations or in terms of position correlations. But these two expansions correspond to different experimental arrangements and cannot be combined into a single context-independent attribution of both position and momentum to the electron at the same time $t$.

%The correlations may be written schematically as
%
%\begin{equation*}
%
%x_E-x_P=0,
%
%\end{equation*}
%
%for the image-plane arrangement, where $x_E$ and $x_P$ are coordinates referred back to the point of collision, and as
%
%\begin{equation*}
%
%p_E-p_P=0,
%
%\end{equation*}
%
%for the focal-plane arrangement, if $p_E$ and $p_P$ denote corresponding momentum transfers with the chosen sign convention. Equivalently, in the usual conservation convention, the latter condition can be written as
%
%\begin{equation*}
%
%\Delta p_E+\Delta p_P=0.
%
%\end{equation*}
%


%% Popper, somewhat prone to paranoia, felt that \vW\ had not taken him seriously and reacted rather aggressively, emphasizing two points: \begin{inparaenum} \item in the Note he did not \emph{prove} the legitimacy of {calculating the past of the electron}, but merely \emph{assumed} it, citing Heisenberg's own remark on p.~15 of his book as justification; \item von {Weizsäcker} had refuted only one {method} of \emph{momentum measurement}---that by scattering. Popper maintained instead that \emph{momentum selection} of an aggregate of electrons \origg{Teilchenmenge} by means of filters could determine momentum \mom \emph{without} disturbing the \pos-coordinate, thus allowing for a non-predictive retrodiction before $X_1$. \end{inparaenum} Popper wrote to Heisenberg:
%
%%\qt{You will now ask why I did not elaborate further in the note on what I am writing here. The reason is that I assumed that the so-called \s{non-predictive measurement} (momentum selection followed by a position measurement) and the possibility of calculating the \s{past} prior to the momentum measurement were already known. This assumption was based on a remark of yours (\textit{Phys. Prinzipien}, p. 15) concerning case (b) and the \s{calculation of the past of the electron}. In my note, however, I had to be brief and restrict myself to what was essential, that is, above all, to what was new. Only from Mr. Weizsäcker's note did I realize that at least he was not familiar with this case}{Sie werden nun fragen, warum ich des, was ich hier schreibe, in der Note nicht näher ausgeführt habe. Der Grund ist, der, dass ich annahm, dass die fragliche \s{nichtprognostische Messung} (Impulsaussonderung mit nachfolgender Ortsmessung) und die Möglichkeit, die \s{Vergangenheit} vor der Impulsmessung zu berechnen, bekannt ist. Diese Annahme gründete sich auf eine Bemerkung von Ihnen (Phys. Prinzipien, S. 15) über den Fall (b) und die \s{Rechnung über die Vergangenheit des Elektrons}. In meiner Note musste ich mich aber kurz fassen und nur das Wichtigste bringen, also vor allem das Neue. Erst aus Herrn Weizsäckers Note bemerkte ich, dass zumindest ihm dieser Fall nicht bekannt ist, und ich muss nun annehmen, dass er entweder auch Ihnen nicht bekannt war (und ich aus der zitierten Stelle Ihres Buches zuviel herausgelesen habe), oder dass Sie, was ich sehr begreiflich finden würde, meine Note nicht recht ernst genommen haben, so dass Ihnen dieser Zusammenhang entgangen ist}[\letterKPAp{Popper}{Heisenberg}{26}{11}{1934}[305-52]] 
%%
%%Since he mentions p.~15 of the German edition of 
%
%
%
%%One might object, Popper points out, that we do not really know in detail what a light filter does, and that its operation might itself disturb the light quantum's state. However, Popper argues, a filter produces not only light but also sharp images \origg{Bilder}---thus, it does not, in any case, disturb the directions of the quanta. If the filter did disturb the light quanta, one would have to assume some mechanism that restores their direction after passage: either (1) the quanta pass through the filter but are unpredictably displaced along their paths, with their previous direction suddenly reestablished \origg{Zurückversetzung auf der Bahn}, or (2) the filtering process involves absorption and directed re-emission \origg{gerichtete Emission} in the same direction. \autocite[407]{Popper1934b} Yet both options appear problematic: the first contradicts the continuity of paths still required by quantum theory; the second assumes a kind of \s{directed fluorescence}\footnote{\s{Fluorescence} is the emission of electromagnetic radiation by atoms or molecules after absorbing radiation, occurring when they return to a lower energy state. The emission is generally spontaneous and isotropic, not directed} not known to occur. \autocite[407]{Popper1934b} Popper therefore concluded that the most plausible hypothesis is that, also in the case of the filter, the quanta with the \s{right} momentum are allowed to pass through unchanged. 
%
%\subsection{Popper's Reply and \VW's Counter-Reply}
%\label{sub:heisobjection}

% D: The slit system (diaphragm) of width $d$. This restricts the transverse position of the incoming particle beam.
%- A and B: The plates of a capacitor (or deflecting field) that bend the particle trajectory.
%- C: The curved trajectory of the particle inside the field.
%- E-F: The vertical displacement $y$ at the detector plane (the "screen" at the right).
%- $L$ : The free flight distance $L$ (the length of the field region).






%Instead, it says: even if the momentum before and after the measuring interaction are both known, the measurement context still does not license a trajectory, because the transition between the two momenta cannot be temporally localized. Hence the later position cannot be propagated backward through the measuring interaction.

%An exact momentum measurement does not by itself yield a trajectory. It fixes the momentum within one experimental context, but leaves the conjugate position insufficiently determined. Conversely, a later position measurement can fix a position only within a different context, and cannot be combined with the preceding momentum measurement to reconstruct a single continuous path.

%A momentum measurement may determine the momentum relevant to a given experimental context, but precisely for that reason it does not also determine the corresponding position history. What is lost is not necessarily knowledge of the momentum before or after the measuring interaction, but the possibility of combining this momentum information with a continuous trajectory through the interaction. The measurement context that licenses the ascription of momentum leaves the conjugate position insufficiently determined; hence it cannot be used to reconstruct the particle's path.


%Hermann initially formulated the objection in terms of disturbance: an exact measurement of a momentum component cannot be interpreted as leaving that same component uncontrollably undisturbed; otherwise it would be unclear how the conjugate position component could be affected. On this view, a momentum measurement can determine either the momentum before or the momentum after the measurement, but not both. \vW's reply to Hermann was more precise. The decisive point is not that the two momenta cannot in principle both be known. Rather, the more exact the momentum measurement, the longer the measuring interaction lasts; and within this interval one cannot determine the instant at which the collision between apparatus and object occurred. Thus one cannot know how long the particle possessed the momentum before the collision and how long it possessed the momentum after it. This temporal indeterminacy is enough to prevent a later position from being traced back through the measurement process to an earlier position with arbitrary precision.


%One can indeed, know the momentum before by reconstructing the collision with the measuring apparatus. If the incoming and outgoing photnon momentum are known, conservation of momentum lets one infer the electrons momentum before and after the measuring collision. But to know the photon momentum very precisely, the probe must be close to a monochromatic wave. A monochromatic wave is extended in time, so the collision cannot be assigned to an exact instant. 







% The central problem is not simply that one cannot know the momentum before and after the measurement. Rather, a sufficiently precise momentum measurement necessarily takes a correspondingly long time. During this interval, one cannot know when the collision between measuring apparatus and measured object occurred. Consequently, one cannot know how long the object had the \s{momentum before the collision} and how long it had the \s{momentum after the collision.} This temporal indeterminacy is enough to blur any position knowledge that one might try to trace backward through the measurement process.


%
%A \s{point counter} (\german{Spitzenzähler}) is essentially a Geiger-type counter in which the sensitive electrode is a fine metal tip $S$, maintained at a high potential difference relative to the counter wall $W$. An ionizing particle passing sufficiently close to the tip produces primary ions in the gas; the strong electric field near $S$ then amplifies this ionization into a discharge. 

	

%Weizsäcker also denies that “momentum selection” differs in principle from “momentum measurement.” Usually, consdiered one Popper's achivenet, to disticntu between preparation and measure. However, that a more orthodocs that any measuremetn, is a preparatio. Howeer, is wha ... The reason is the temporal symmetry of nature with respect to past and future: selecting particles with a given momentum and measuring their momentum are not physically distinct operations in the sense Popper requires. Still, Weizsäcker is cautious about whether this \s{recipe for refuting Mr. Popper} can be turned into a fully general, compelling argument. In the end, he suggests, one must appeal to the quantum-mechanical formalism as a whole, which Popper appears not to have understood. That that selection, put in auotstate of coresponding operato.


% Any arrangement that would allow one to undercut the uncertainty relations would have to make possible the construction of a wave group with $\frac{\Delta q}{\Delta \lambda}<1$. But such a construction would contradict ordinary wave kinematics. Hence, whenever one analyzes a Popperian arrangement purely in wave-theoretical terms, one should expect to find a contradiction. 





%spatial coordinates correspond directly to scattering angles,


%If you move the object plane infinitely far away ( $d_o \rightarrow \infty$ ), the light rays reaching the lens are parallel.

%When you move the object plane closer to the lens (a finite $d_o$ ), the light rays from the object reach the lens at diverging angles.

%To switch to the image plane of a nearby object, you must increase the physical distance between the lens and the sensor.

%Imagine the object plane is located at an infinite distance from the lens.

%As you move the subject closer to or further from the lens, the object plane moves accordingly.

%$$\frac{1}{f}=\frac{1}{d_o}+\frac{1}{d_i}$$

%n short, the focal plane is where the lens focuses an object located at infinity, serving as a baseline property of the lens itself. The image plane is where the lens focuses the specific object you are currently looking at, changing position based on how far away that object is.

%The focal plane is where parallel rays are focused; the image plane is where rays from an object point are brought back together.

%The same scattering event can be inserted into two different experimental preparations. If the plate is placed in the image plane, the optical system prepares a point-to-point correlation P- P, so the mark is interpretable as position information. If the plate is placed in the focal plane the optical system prepares a direction-to-point correlation k I P', so the mark is interpretable as momentum information. The mark itself is therefore not a self-standing datum; its physical meaning is fixed by the preparation of the optical arrangement.


%In the image-plane arrangement, the apparatus prepares the registration so that rays originating from the same object point are recombined at the same image point:


%He then turns to Hermann’s second point and says that he can agree with her account almost completely. His only reservation comes from his limited knowledge of Popper’s book and from the difficulty he has had in determining exactly what Popper misunderstood. He recalls that Popper had earlier treated his \s{non-prognostic measurements} as if they were a generally known and accepted matter. Weizsäcker notes that Popper’s book had already been printed before Popper saw his reply; Popper intended to publish responses in a journal but had not yet done so because his claims were being refuted immediately in correspondence. Still, Weizsäcker grants that Popper is reasonable enough to abandon views once they are clearly refuted. For that reason, he had thought Popper might have something interesting to say in epistemology, even if his physics was weak.
%
%On the question whether wave-particle dualism can be demonstrated in processes involving individual particles, Weizsäcker is cautious. He thinks that, for individual particles, the dualism can hardly be shown except by repeating experiments many times. He cannot think of a single, one-off experiment that would count as a proof of the wave picture. The reason, he suggests, is that all measurements ultimately end in the determination of a position at a time, that is, a spatiotemporal coincidence. This always permits a corpuscular mode of expression for the individual case.
%
%His example is an electron with known momentum incident on a grating of thin wires. According to the corpuscular picture, the electron should be scattered with equal probability in arbitrary directions. According to the wave picture, however, diffraction occurs, so that certain deflection directions are privileged. If the electron is later found in one of these directions, one may say that without wave properties it would have been too improbable for it to be deflected precisely that way. But this reasoning already involves probability, and probability becomes meaningful only through many trials. Thus, in the single case, the wave aspect is not directly demonstrable in the same way it is in a repeated series of experiments.

%Within each context, wave-particle duality regulates the reciprocal limitation, so that e.g. a sharp momentum the position must completely undetermined. 

%Die zu verschiedenen Beobachtungszusammen-
%hängen gehörenden Bestimmungsstücke werden also in einer einzigen Beschrei-
%bung des physikalischen Systems vereinigt. Da nach dem quantenmechanischen
%Formalismus jede physikalische Zustandsbeschreibung nur relativ zum jeweiligen
%Beobachtungszusammenhang gilt, so ist diese Vereinigung der zu verschiede-
%nen Beobachtungszusammenhängen gehörenden Bestimmungsstücke in einer
%Beschreibung physikalisch bedeutungslos



%1. Die Messung einer Impulskomponente, bei deren Interpretation nicht mit einer  unkontrollierbaren Störung eben dieser Impulskomponente gerechnet werden müsste, ist mit dem Formalismus unvereinbar.

%For otherwise it would be impossible to see—as P. also assumes—that a disturbance of the corresponding position component could occur. Thus, if one had determined the system before the measurement with maximal accuracy with respect to these position and momentum components, then the measurement would immediately exceed the limits imposed by the uncertainty relations.


%Popper’s argument requires such a measurement because the selected momentum must not merely characterize the particle between the momentum measurement and the subsequent position measurement; it must also determine the particle’s path before the momentum measurement, so that the earlier collision and the trajectory of the correlated particle can be reconstructed with a precision exceeding the uncertainty relations. Hermann therefore asks whether the impossibility of such a disturbance-free momentum measurement can be shown from the experimental methods themselves, and more generally from the requirements of the dualism of wave and particle descriptions, which must enter into the interpretation of every measurement.

%Die bei der Interpretation einer Messung auftretenden quantenmecha- nisch zulässigen und notwendigen Rückschlüsse unterscheiden sich von diesen Rechnungen gerade dadurch, dass sie den Dualismus von Wellen- und Partikelbild berücksichtigen

%In the relevant optical arrangements, the apparatus can be set up either to correlate a mark on the plate with the point from which the light was scattered, or to correlate it with the direction in which the scattered light propagated. The first arrangement yields position information; the second yields momentum information. What it cannot do is provide, in one and the same context, both a sharply defined point of origin and a sharply defined direction of propagation.

%Such calculations appear to overcome the uncertainty relations because they combine elements fixed in different observations, for example two position measurements or a position and a momentum measurement, into a single classical description of a past trajectory. But the determinations so combined belong to different observational contexts and cannot be fused into one physically meaningful state description. Referring briefly, to the \vW's $\gamma$ ray microscope example that would play a central role in her book, that Popper only used a momentum measuring apparatus, that to does not allow of the exact determination of position.



%A quantum measurement cannot be interpreted in purely corpuscular terms, as if the apparatus merely registered the pre-existing path of a particle. Rather, the measurement correlates a wave-like process with a corpuscular event. This correlation can be established in different ways, depending on how the apparatus is used. In one arrangement, the wave process is correlated with a point of origin and thereby with the position of the event; in another, it is correlated with a direction of propagation and thereby with the momentum exchange. The choice between these arrangements defines different observational contexts. Within each context, the wave and particle descriptions then mutually restrict one another: a context that yields a relatively sharp position leaves the corresponding momentum indeterminate, whereas a context that yields a relatively sharp momentum leaves the corresponding position indeterminate. Popper’s error is to treat the resulting determinations as if they could be detached from their contexts and combined into a single trajectory


%Hermann’s criticism can be reconstructed in three steps. 
%
%First, Popper’s immediate error is to project the result of a momentum measurement backward before the measurement itself. He treats the momentum selected at the first apparatus as if it also determined the particle’s earlier momentum, and hence as if it allowed him to reconstruct the path leading back to the collision. At this point, Hermann since that disturebs the the 
%
%In third point, Hermann’s criticism goes deeper. Even the familiar retrospective trajectory calculations based on successive measurements do not yield a physically meaningful path. They combine determinations belonging to different observational contexts, for example a momentum determination from one context and a position determination from another, and fuse them into a single classical trajectory. 
%
%The second point, Hermann traces both errors back to Popper’s misunderstanding of wave-particle duality. 

%Because he treats the wave function as a statistical description of ensembles of particles, rather than as part of the context-relative description of individual systems, Popper fails to see that each observational context is constituted by a specific coordination of wave and particle descriptions. He therefore treats position and momentum determinations as if they revealed different aspects of one underlying corpuscular path, instead of recognizing that they are meaningful only within the contexts in which the wave and particle pictures mutually restrict one another.

%Hermann therefore rejects the idea that the physical reality of such past trajectories is a mere matter of taste: the reconstructed path is physically meaningless insofar as it neglects the dependence of every physical determination on its \german{Beobachtungszusammenhang}. 





%The contrast can be reconstructed in three steps. First, Popper treats the \psi-function as a wave-like statistical description of an ensemble of particles. The particles themselves are still assumed to possess determinate corpuscular paths; the wave function merely gives the probability distribution for such particles within an ensemble. Second, Hermann objects that wave-particle duality applies already to the description of individual physical systems, not merely to ensembles. A quantum measurement therefore cannot be interpreted in purely corpuscular terms, as if the apparatus merely registered the pre-existing path of a particle. Rather, the measurement correlates a wave-like process with a corpuscular event. This correlation can be established in different ways, depending on how the apparatus is used. In one arrangement, the wave process is correlated with a point of origin and thereby with the position of the event; in another, it is correlated with a direction of propagation and thereby with the momentum exchange. The choice between these arrangements defines different observational contexts. Third, within each context, the wave and particle descriptions mutually restrict one another: a context that yields a relatively sharp position leaves the corresponding momentum indeterminate, whereas a context that yields a relatively sharp momentum leaves the corresponding position indeterminate. 

%Popper’s error is to detach determinations from the contexts in which they are meaningful and to combine them into a single position-cum-momentum trajectory.



%A quantum measurement cannot be interpreted in purely corpuscular terms, as if the apparatus merely registered the pre-existing path of a particle. Rather, the measurement correlates a wave-like process with a corpuscular event. This correlation can be established in different ways, depending on how the apparatus is used. In one arrangement, the wave process is correlated with a point of origin and thereby with the position of the event; in another, it is correlated with a direction of propagation and thereby with the momentum exchange. The choice between these arrangements defines different observational contexts. Within each context, the wave and particle descriptions then mutually restrict one another: a context that yields a relatively sharp position leaves the corresponding momentum indeterminate, whereas a context that yields a relatively sharp momentum leaves the corresponding position indeterminate. Popper’s error is to treat the resulting determinations as if they could be detached from their contexts and combined into a single trajectory.


%The electron-photon interaction cannot be described by a purely corpuscular collision picture alone. The propagation of the scattered light through the microscope must be described in the wave picture, while the interaction between the light quantum and the electron is described in the corpuscular picture





%Hermann’s structure can be reconstructed as follows:
%
%1. Wave-particle duality first makes the measurement problem take this form.
%    The electron-photon interaction cannot be described by a purely corpuscular collision picture alone. The propagation of the scattered light through the microscope must be described in the wave picture, while the interaction between the light quantum and the electron is described in the corpuscular picture. The apparatus is therefore needed not merely to “see” the collision, but to establish how the wave-optical propagation is correlated with the corpuscular scattering event.
%2. This correlation generates the inter-contextual alternative.
%    Once the microscope is treated as an optical apparatus, there are two relevant ways of using it. Observation in the image plane correlates the observed wave with a point of origin; in corpuscular terms, this gives the collision point, hence the electron’s position. Observation in the focal plane correlates the observed wave with a direction of propagation; in corpuscular terms, this gives the photon’s momentum direction and hence the momentum exchange. Thus the image-plane/focal-plane alternative is not simply external to wave-particle duality: it arises from the attempt to connect wave-optical propagation with corpuscular collision.
%3. Within each context, wave-particle duality regulates the reciprocal limitation.
%    In the image-plane context, the wave description supports localization at a point of origin, but the corresponding momentum exchange is blurred. In the focal-plane context, the wave description supports a sharp direction of propagation, but the corresponding collision point is blurred. Thus the two pictures do not merely provide two descriptions of the same underlying trajectory; they mutually restrict one another within the selected observational context.

%P. stützt sich bei seinen Überlegungen begreiflicher Weise stark auf die bekannten Bahnberechnungen, die – etwa auf Grund zweier Ortsmessungen oder einer Orts-und einer Impulsmessung – für einen vergangenen Zeitabschnitt den Zustand eines Systems mit einer die Unbestimmtheitsrelationen übersteigenden Genauigkeit angeben. Im Anschluss an den Popperschen Missbrauch mit diesen Rechnungen
%habe ich mir eingehender überlegt, wie eigentlich der Schein einer Überwindung der Unbestimmtheitsrelationen hier entstehen konnte. Ich meine, dadurch, dass diese Rechnungen den Dualismus von Wellen- und Partikelbild nicht beachten sich auf die anschaulichen Merkmale nur eines dieser Bilder beschränken – in den bekannten Fällen: des Korpuskelbildes – und, da jede einzelne Messung nur unter Berücksichtigung beider Bilder interpretiert werden kann und daher keine von ihnen unbeschränkt anzuwenden gestattet, diese Merkmale durch zwei verschiedene Beobachtungen festlegen. Die zu verschiedenen Beobachtungszusammenhängen gehörenden Bestimmungsstücke werden also in einer einzigen Beschrei-
%bung des physikalischen Systems vereinigt. Da nach dem quantenmechanischen
%Formalismus jede physikalische Zustandsbeschreibung nur relativ zum jeweiligen
%Beobachtungszusammenhang gilt, so ist diese Vereinigung der zu verschiede-
%nen Beobachtungszusammenhängen gehörenden Bestimmungsstücke in einer
%Beschreibung physikalisch bedeutungslos


%In the microscope example, observation in the image plane yields a relatively sharp collision point but an indeterminate momentum exchange, whereas observation in the focal plane yields a relatively sharp momentum exchange but an indeterminate collision point. 

%Legitimate quantum-mechanical inferences differ precisely in that a single observation fixes how wave and particle descriptions mutually restrict one another. In the microscope example, observation in the image plane yields a relatively sharp collision point but an indeterminate momentum exchange, whereas observation in the focal plane yields a relatively sharp momentum exchange but an indeterminate collision point. The mutual regulation of the two descriptions preserves the observational context and, with it, the uncertainty relations. Popper’s trajectory reconstruction therefore becomes physically meaningless in the measure in which it exceeds those relations.
\end{itemize}

%Hermann’s criticism of Popper is difficult to reconstruct because her own language is not fully differentiated. She begins from the dualism of wave and particle descriptions and uses it to cover two distinct issues. The first is the choice between incompatible observational contexts: the apparatus may be arranged as a position-measuring device, as in the image plane, or as a momentum-measuring device, as in the focal plane. The second is the reciprocal limitation within each context: once a context has been fixed, wave-particle duality explains why a sharp determination of momentum precludes the corresponding localization, and why a sharp localization precludes the corresponding momentum determination. 

%Hermann’s criticism of Popper is difficult to reconstruct because her own language is not fully differentiated, and the wave particle duality seems to be at the basis of two different issues.


%Popper’s trajectory reconstruction fails because it ignores both levels. He treats the wave function as a statistical device for ensembles of particles with underlying paths, and therefore combines position and momentum determinations that are meaningful only within different observational contexts.

%So the first level already involves wave-particle duality in this sense: optical wave information is translated into corpuscular information about the electron. The image plane gives access to where the photon was scattered, hence to the electron’s position at the collision; the focal plane gives access to the direction of the scattered wave, hence to the photon’s momentum and thereby to the electron’s momentum exchange. But there is then a second, stricter sense of duality within each context. Once the apparatus is used in the image plane, the wave picture gives a spherical wave from a relatively sharp point, while the particle picture gives a collision at a relatively sharp position but with blurred momentum exchange. Once the apparatus is used in the focal plane, the wave picture gives a parallel bundle with a relatively sharp direction, while the particle picture gives a relatively sharp momentum exchange but an indeterminate place of collision.



%The image-plane/focal-plane alternative is not simply external to wave-particle duality. As an optical matter, the image plane and the focal plane offer two different ways of using the microscope: the image plane correlates the observed image with a point of origin of the scattered wave, whereas the focal plane correlates the observation with a direction of propagation of that wave. In the quantum-mechanical interpretation, these wave-optical determinations correspond respectively to corpuscular determinations: the point of origin of the scattered wave corresponds to the position of the collision, while the direction of propagation corresponds to the momentum of the scattered photon, and hence to the momentum exchange with the electron. In this sense, the choice between image plane and focal plane is already an application of wave-particle duality.
%
%There is, however, a second level of duality within each observational context. If the apparatus is used in the image plane, the wave description supports the idea of a spherical wave originating from a relatively well-defined point, and the particle description accordingly yields a collision with relatively sharp position but indeterminate momentum exchange. If the apparatus is used in the focal plane, the wave description supports the idea of a parallel bundle with a relatively well-defined direction, and the particle description accordingly yields a relatively sharp momentum exchange but an indeterminate collision point. Thus the first level concerns the optical alternative between two possible contexts, while the second concerns the reciprocal restriction of wave and particle descriptions within whichever context is chosen.

%First, inter-contextual duality: the optical distinction between image plane and focal plane maps wave-origin information onto particle-position information, or wave-direction information onto particle-momentum information.
%
%Second, intra-contextual duality: within the chosen context, the wave and particle descriptions mutually restrict one another, producing the relevant uncertainty trade-off.

%Hermann tends to describe both issues in the same language of wave-particle duality. She therefore presents Popper’s statistical interpretation of the wave function as a failure to grasp wave-particle duality, while what is at stake is more precisely a twofold failure: Popper fails to recognize both the context-dependence of quantum-mechanical state descriptions and the reciprocal trade-off imposed within each context. Since he treats the wave function as a statistical device for ensembles of particles with underlying trajectories, he assumes that determinations obtained in different contexts can be combined into a single position-cum-momentum path. Hermann’s central objection is that this combination is physically meaningless.

%Hermann’s formulation is not simply mistaken. The alternative between image plane and focal plane is not itself identical with wave-particle duality, but it is also not independent of it. As an optical matter, the image plane and the focal plane offer two different ways of using the microscope: the former correlates the image with a point of origin, the latter with a direction of propagation. In the (\gamma)-ray microscope, however, this optical alternative becomes a quantum-mechanical alternative only because the scattering event must also be described corpuscularly, as a collision between a light quantum and an electron. Observation in the image plane therefore corresponds to a relatively sharp determination of the collision point, while observation in the focal plane corresponds to a relatively sharp determination of the momentum exchange. Thus the choice of observational context is grounded in the wave-optical structure of the apparatus, but its significance as a choice between position and momentum measurements depends on the simultaneous use of the particle picture. Hermann’s language is compressed rather than simply confused: wave-particle duality does not by itself constitute the context alternative, but it explains why the optical alternative becomes an alternative between incompatible quantum-mechanical determinations.

%The logical structure of Hermann’s criticism can be reconstructed as follows. First, Popper reads the wave function statistically, as a probability distribution over particles that are still assumed to possess classical paths, that is, determinate position-cum-momentum histories. The wave function therefore characterizes an ensemble rather than describing an individual system relative to a definite observational context. Second, this statistical reading leads Popper to miss the context-dependence of quantum-mechanical state descriptions. If the individual particle already has a context-independent trajectory, then different measurements may be treated as revealing different aspects of that same path; a position determination and a momentum determination can accordingly be combined in a retrospective trajectory calculation. Third, Popper fails to apply wave-particle duality to each individual system within each measurement context. For Hermann, a momentum context permits a sharp momentum determination only through a wave-like description that precludes the corresponding localization, while a position context permits localization only at the cost of the corresponding momentum determination. Popper’s error is therefore that he treats the wave function as an ensemble calculus over underlying trajectories, and thereby combines determinations that are meaningful only within distinct observational contexts.

%Recent reconstructions of \vW’s example make this distinction clearer: the total photon-electron state may be represented in different bases, and the relative state assigned to the electron depends on the measurement context. 

%Popper’s statistical reading of the wave function explains why he fails to see the contextual character of quantum-mechanical state descriptions. If the wave function is only a mathematical device for describing ensembles, then the uncertainty relations restrict only the homogeneity of such ensembles; they do not prevent an individual particle from having simultaneously determinate position and momentum. On this assumption, separate measurements may be treated as revealing different aspects of one and the same underlying trajectory. Hermann rejects precisely this move. For her, the wave function describes an individual system relative to a definite observational context. A momentum measurement and a position measurement therefore do not disclose two independently existing parts of a single classical path; they define different physical determinations within different contexts. Popper’s trajectory reconstruction is illegitimate because it combines these context-relative determinations as if they belonged to one context-independent description.


%Hermann’s criticism of Popper is still somewhat disorganized at this point, because two issues remain entangled. The first is the alternative between incompatible observational contexts. In the image plane, the apparatus functions as a position-measuring arrangement; in the focal plane, it functions as a momentum-measuring arrangement. The second issue is the coordinated use of wave and particle descriptions within each context. Once an observational context has been fixed, the measurement must still be interpreted through the mutual restriction of the wave and particle pictures. The more precise point, therefore, is not simply that Popper ignores wave-particle duality, but that he combines determinations belonging to different observational contexts, while wave-particle duality explains why each context imposes its own reciprocal limitations.

%Popper’s deeper mistake can then be stated more sharply. He chooses, in effect, a momentum context: the apparatus is supposed to select particles with a sharply determined momentum component. But within such a context the wave description precludes the positional determination needed for a trajectory. A sharp momentum selection corresponds to a plane-wave-like state, and hence to the loss of localization in the conjugate coordinate. Popper therefore cannot use the momentum selected in this context to reconstruct a determinate path through the apparatus and back to the collision point. The difficulty is not that momentum cannot be measured sharply at all, but that a sharp momentum measurement cannot simultaneously support the position determinations required for a classical trajectory.

%This also explains why Popper misses the point. He treats the wave function as a statistical or mathematical device for describing ensembles of particles, while the particles themselves are still imagined as possessing determinate classical trajectories. He therefore takes the uncertainty relations to limit only the homogeneity of ensembles, not the applicability of classical concepts within an individual measurement context. In Hermann’s terms, the order of explanation should be reversed: Popper’s mistaken statistical reading of the wave function explains why he fails to see the force of wave-particle duality, and this in turn explains why he illegitimately combines determinations from incompatible observational contexts.

%Popper’s statistical reading of the wave function explains why he fails to see the contextual character of quantum-mechanical state descriptions. If the wave function is only a mathematical device for describing ensembles, then the uncertainty relations restrict only the homogeneity of such ensembles; they do not prevent an individual particle from having simultaneously determinate position and momentum. On this assumption, separate measurements may be treated as revealing different aspects of one and the same underlying trajectory. Hermann rejects precisely this move. For her, the wave function describes an individual system relative to a definite observational context. A momentum measurement and a position measurement therefore do not disclose two independently existing parts of a single classical path; they define different physical determinations within different contexts. Popper’s trajectory reconstruction is illegitimate because it combines these context-relative determinations as if they belonged to one context-independent description.



%Hermann's use of wave-particle duality involves two levels that should be kept distinct. At the first level, there is the choice between different observational contexts. In \vW's microscope example, observation in the image plane defines a position context: the collision point is relatively sharp, while the momentum exchange is blurred. Observation in the focal plane defines a momentum context: the momentum exchange is relatively sharp, while the collision point is blurred. At the second level, each of these contexts must itself be interpreted through a coordination of the wave and particle pictures. The propagation of light through the microscope is described in the wave picture, while the collision between light quantum and electron is described in the particle picture; the aperture angle $\epsilon$ then expresses how the wave-optical conditions restrict the corpuscular inference. 



%Popper's mistake is to ignore this contextual organization: he combines the position determination belonging to one context with the momentum determination belonging to another, and thereby constructs a trajectory that is not licensed by any single quantum-mechanical observational context.


%Strictly speaking, the image-plane/focal-plane alternative should not itself be called a form of wave-particle duality. It is rather an alternative between incompatible observational contexts: in the image plane the apparatus functions as a position-measuring arrangement, whereas in the focal plane it functions as a momentum-measuring arrangement. Wave-particle duality enters at a different level. Once an observational context has been fixed, the measurement must be interpreted through a coordinated use of the wave and particle pictures. In the image-plane context, this coordination permits a relatively sharp localization of the collision point but leaves the momentum exchange indeterminate; in the focal-plane context, it permits a relatively sharp determination of the momentum exchange but leaves the collision point indeterminate. Hermann's wording can therefore be misleading, since she describes the incompatibility of measurement contexts in the language of wave-particle duality. The more precise point is that Popper combines determinations belonging to different contexts, while the duality of wave and particle descriptions explains why each context imposes its own reciprocal limitations.

%Hermann’s three objections are not yet fully coherently ordered. At first she locates Popper’s mistake in the assumption of a momentum measurement that does not disturb the measured momentum component itself. In this formulation, the difficulty seems to be that Popper wants to determine the momentum of the particle both before and after the measurement. But this way of putting the point is somewhat clumsy. The deeper difficulty is not simply that the momentum cannot be determined; a measurement context may indeed determine momentum sharply. The problem is rather that such a determination excludes the positional information required for a trajectory. Once the experiment is arranged as a momentum measurement, the corresponding position determination is no longer available in the same observational context, and conversely. Popper’s mistake is to combine determinations belonging to different contexts as if they formed a single classical description of the particle’s path.
%
%The second point explains why Popper fails to see this. According to Hermann, he is misled by the statistical interpretation of the wave function. He treats the wave function as a mathematical device for characterizing ensembles of particles, rather than as a state description of an individual system relative to a definite observational context. On this reading, the electrons remain, as it were, ordinary particles with definite but perhaps unknown classical properties, while the wave function merely gives probabilities for their distribution in an ensemble. The uncertainty relations then appear to Popper as restrictions on the homogeneity of ensembles, not as limitations on the simultaneous applicability of classical concepts such as position and momentum to individual systems.
%
%The more fundamental point, however, is the third. For Hermann, every measurement context fixes a particular coordination of wave and particle descriptions. If one measures position, the context permits a relatively sharp localization, but only at the cost of the conjugate momentum determination. If one measures momentum, the context permits a relatively sharp momentum determination, but only at the cost of localization. Wave-particle duality does not by itself determine which context is chosen; rather, once a context is fixed by the experimental arrangement, duality governs how the wave and particle descriptions mutually restrict one another within that context. Popper’s retrospective trajectory calculations fail because they detach the determinations from their respective contexts and combine them into a single trajectory that no single quantum-mechanical observation can support.
%
%
%
%
%
%Popper thinks in terms of particles plus wave-like probabilities, whereas Hermann wants him to think in terms of context-dependent applications of the wave and particle pictures.



%
%
%
%%Hermann's objection is that Popper's arrangement requires a momentum measurement to have an illegitimate double function. It must determine the electron's momentum between $X_1$ and $X_2$, but it must also determine the electron's momentum before $X_1$, since only then can Popper reconstruct the earlier path of the electron and fix the collision at $S$. 
%
%%Hermann first tentative response is that Popper's arrangemet this is incompatible with the formalism: a measurement of a momentum component cannot be interpreted as one in which no uncontrollable disturbance of that very component has to be taken into account. A precise momentum measurement can therefore determine either the momentum before the measurement or the momentum after the measurement, but not both. Otherwise the measurement would allow one to combine a maximally precise determination of the relevant position component with a maximally precise determination of the corresponding momentum component, and would thereby cross the limits imposed by the uncertainty relations. \q{Stimmt das?}, she asks \vW. 
%
%
%%Popper needs an impulse measurement that does not disturb the very momentum component it is supposed to determine. She now asks whether this can be shown to be impossible not merely by reviewing particular experimental techniques, but by appealing to the general requirements imposed by wave-particle duality. At this point, Hermann hopes to find a response every measurement must be interpreted through the coordinated use of wave and corpuscular descriptions, and this coordination should reveal why Popper’s allegedly non-disturbing momentum measurement is physically inadmissible. Hermann then turns to Popper’s own proposed devices: a filter placed in a light beam, spectral analysis of an electron beam by means of an electric field perpendicular to the direction of propagation, and the point counter. In the case of the electric field she already suspects that a change of momentum is unavoidable, but she does not yet know exactly how to derive the uncontrollability of this change from the dualistic interpretation of the experiment
%
%%The deeper reason of mistake howver, is  
%%
%%Hermann’s point is that Popper is misled by familiar retrospective trajectory calculations. Such calculations appear to determine the state of a system for a past interval with a precision exceeding the uncertainty relations, for example from two position measurements or from one position and one momentum measurement. Hermann explains this appearance by saying that these calculations ignore wave-particle duality. 
%
%%Hermann's criticism is not simply that Popper neglects wave-particle duality, as if duality were an independent rule added to the formalism. Rather, her point is that every measurement context fixes a specific way in which the wave and particle descriptions are to be coordinated. A position measurement defines one observational context: within it the particle aspect can be used to assign a relatively sharp position, while the wave aspect limits the corresponding momentum determination. A momentum measurement defines another context: within it the wave aspect can be used to assign a relatively sharp momentum, while the corresponding position determination is lost. Popper's mistake is to take determinations belonging to different contexts and combine them into a single classical trajectory. For Hermann, this is physically meaningless, because a quantum-mechanical state description is valid only relative to the observational context in which the relevant quantities are defined.
%
%
%They combine determinations belonging to different observational contexts into a single classical description, usually in terms of the corpuscular picture alone. But, according to Hermann, quantum mechanics permits physical state descriptions only relative to a definite observational context. A description that combines elements fixed in different contexts is therefore physically meaningless, and this meaninglessness shows itself in the fact that such a reconstructed trajectory is uncontrollable and cannot serve as the basis for predictions. Legitimate quantum-mechanical inferences differ from these fictitious trajectory calculations because they preserve the unity of the observational context: a single observation fixes how the wave and particle pictures mutually restrict one another. In the microscope example, observation in the image plane gives a relatively sharp collision point but an indeterminate momentum exchange; observation in the focal plane gives a relatively sharp momentum exchange but an indeterminate collision point. The crucial point is that the two pictures regulate one another within one context, thereby preserving the uncertainty relations. Hermann therefore proposes to correct Heisenberg’s remark that attributing physical reality to such retrospective trajectories is merely a matter of taste: these calculations neglect the dependence of every physical determination on its observational context and are, in this sense, physically meaningless.
%


%Hermann argues that Popper’s deeper mistake is to treat the wave function, because of its probabilistic interpretation, as a description only of ensembles of similarly prepared systems rather than as a state description of an individual system relative to a definite observational context. Since Popper understands probability statements in terms of relative frequencies, he also takes the uncertainty relations to apply only to ensembles: they would merely state that no ensemble can be prepared more homogeneously than a pure case. He therefore thinks that individual systems might still, under suitable experimental arrangements, be found in states determined more precisely than the uncertainty relations allow. Hermann rejects this. In her view, the uncertainty relations follow from the experimental necessity of applying both wave and particle descriptions to every atomic process. They therefore limit not merely ensembles, but the very applicability of classical state descriptions to individual systems. The probability interpretation captures only one side of the matter: the side concerning transitions between observational contexts. Complementarily, the wave function also gives an exact quantum-mechanical description of a system relative to the present observational context, namely through the observables for which it is in an eigenstate. Popper’s error, then, is to miss this side of Bohr’s complementarity: within a given context, the system is exactly described, but classical double attributes such as position plus momentum are defined only within the limits fixed by the uncertainty relations.



%The point is methodological: every measurement must be interpreted through the coordinated use of wave and corpuscular descriptions, and this coordination should reveal why Popper’s allegedly non-disturbing momentum measurement is physically inadmissible. Hermann then turns to Popper’s own proposed devices: a filter placed in a light beam, spectral analysis of an electron beam by means of an electric field perpendicular to the direction of propagation, and the point counter. In the case of the electric field she already suspects that a change of momentum is unavoidable, but she does not yet know exactly how to derive the uncontrollability of this change from the dualistic interpretation of the experiment


%Thus the contrast is not simply between electron $A$ and photon $B$, but between a momentum selection with given direction and a full trajectory reconstruction requiring both magnitude and direction. 


%to measure the momentum of electron $A$, but to show that this measurement determines the momentum that $A$ already had before the momentum measurement. Popper needs this retrodictive step because only then can he reconstruct the earlier trajectory of $A$, and from a previous collision infer the trajectory of the correlated particle $B$ with a precision exceeding the uncertainty relations.

%Hermann thus identified immediatly the issue on which the entire hinges, that is the possbility of determining the momentum before the first momentum measuremnt.  Indeed, Popper himself idengied it quite clearlu in appendix VI.

%Hermann's draft formulation: momentum measurement disturbs momentum itself, therefore one cannot know both pre- and post-measurement momentum. A precise momentum measurement may determine either the momentum before the measurement or the momentum after the measurement, but not both at once. The reason is structural. If a measurement could determine a momentum component without any uncontrollable disturbance of that component, then there would be no corresponding reason for the conjugate position component to become disturbed. But if both position and momentum had previously been fixed with maximal accuracy, such a measurement would immediately violate the uncertainty relations. 



Hermann therefore asks whether one can show, not merely by discussing individual measuring devices, but to find a principled answer, that there can be no momentum measurement that leaves the measured momentum component undisturbed. Her proposed route is to appeal to the wave-particle duality involved in every measurement interpretation. Any concrete measurement must combine the wave and corpuscular pictures; and the way these pictures limit one another should explain why a momentum measurement necessarily involves an uncontrollable disturbance, or at least why it cannot determine both the pre-measurement and post-measurement momentum. She then applies this question to Popper’s examples. He also claims that an electron beam can be analyzed by means of an electric field normal to the beam direction. Hermann finds the second case clearer: in an electric field, a momentum change seems unavoidable. What she lacks is a precise account of how the uncontrollability of that change follows from the wave-particle dualism. In the case of a light filter, Popper claims that a light ray’s momentum component can be measured without disturbing it by placing a filter in the beam.

In sum, Hermann’s question is not merely technical but foundational: she wants to know how to prove, from the quantum-mechanical structure of measurement itself, that Popper’s allegedly non-disturbing momentum measurements are impossible. Her suspicion is that Popper’s proposal fails because it assumes that one can isolate a momentum value without paying the complementary price imposed by the wave-particle dualism and the uncertainty relations.

2. Since Hermann's could read the \LdF, she realize the deeper assumption behind Popper's stance than \vW.  Hermann identifies the deeper source of Popper’s physical mistakes in his interpretation of the wave function. In her view, Popper is misled by the probabilistic interpretation into treating wave functions not as descriptions of individual physical systems, but only as characterizations of ensembles of systems. Since probability statements can be understood as statements about relative frequencies in a series of cases, Popper concludes that the wave function concerns only such series. From this, Popper also infers that the uncertainty relations apply only to ensembles. On this reading, the uncertainty relations say merely that no ensemble can be prepared that is more homogeneous than a pure quantum case. They would not exclude the possibility that, under a suitable experimental arrangement, one might encounter individual systems whose state is determined more precisely than the uncertainty relations allow.


Popper interprets the uncertainty relations as scattering relations: they govern ensembles of particles whose statistical distribution is wave-like, while the individual particles may still be conceived corpuscularly. Hermann objects that this misses their real source. The relations do not follow merely from uncontrollable scattering interactions, but from wave–particle duality itself: each individual electron or photon can be treated, depending on the experimental arrangement, under either the wave or the particle description. The uncertainty relations express the limits within which these two descriptions can be jointly applied. To select a definite momentum is to select a plane wave; but a plane wave is not localized. Hence an electron with perfectly well-defined momentum has a completely undetermined position. 


%In the ideal limit, exact momentum selection corresponds to a plane wave and therefore to complete spatial indeterminacy. In any real apparatus, approximate momentum selection corresponds to an extended wave packet: the narrower the momentum spread, the greater the indeterminacy of position.


%%He imagines an apparatus that allows only electrons with a certain momentum, or a sufficiently narrow momentum range, to pass. In a corpuscular picture, this looks straightforward: electrons are little particles with definite momenta, and the apparatus simply filters out all particles except those whose momentum has the desired value. This is the sense in which Popper is working with a particle picture.
%
%%Popper treats the uncertainty relations as scattering relations: the wave-like element belongs primarily to the statistical distribution of an ensemble, while the individual particle may still be conceived corpuscularly. Hermann objects that this misses the structure of the measurement itself. In each individual measurement, the electron and photon must be treated, according to the stage of the experiment, both as particles and as waves. The uncertainty relations express the conditions under which these two descriptions can be consistently coordinated in a given experimental arrangement.
%
%%But Hermann would say that this is incomplete. To make the momentum selection physically precise, the electron beam must be constrained by an actual experimental arrangement: slits, diaphragms, fields, apertures, or some equivalent selection device. The moment one asks how the selection is implemented, one has to treat the electron also as a wave. A sharply selected momentum corresponds to a wave with a well-defined wavelength and direction. But producing such a momentum selection destroys localization.
%
%\todo{uncertainty relations}

Hermann regards this as the decisive error. The uncertainty relations do not arise merely from the statistical interpretation of the wave function. Rather, following Heisenberg, they arise from the dualism experiments and from the necessity of applying wave-like features to every atomic process and corpuscular features to every wave process. This means that the limits expressed by the uncertainty relations constrain not only ensembles but also the possible state-description of an individual physical system. Popper’s momentum filter can be understood, in the first instance, in corpuscular terms: the apparatus is supposed to select only those electrons whose momentum has a specified value. Hermann’s objection is that this picture becomes untenable as soon as one asks what such a precise selection physically amounts to. To select a sharply defined momentum is, in wave-mechanical terms, to select a plane wave, that is, a state with a definite wavelength and direction of propagation. But a plane wave is completely delocalized. Thus the very conditions that make the electron’s momentum well defined make its position wholly indeterminate. The apparatus cannot therefore be understood as merely sorting localized particles according to pre-existing momenta; the wave description is already implicit in the physical possibility of the momentum selection itself.

The uncertainty relation is the quantitative expression of the incompatibility between the two possible measurement contexts. It says how much indeterminacy in $q$ is necessarily entailed by an arrangement that makes $p$ determinate, and conversely how much indeterminacy in $p$ is entailed by an arrangement that makes $q$ determinate. 




%Popper’s experiment is designed as a momentum measurement. He assumes that the apparatus can filter electrons according to their momentum, so that only electrons with a sufficiently well-defined momentum pass through. What this description overlooks, from Hermann’s point of view, is that the electron cannot be treated merely as a corpuscle selected by the apparatus. The more precisely its momentum is determined, the more the corresponding wave must approximate a plane wave; and the more it approximates a plane wave, the less determinate its position becomes. A momentum-measuring apparatus therefore does not also allow the electron’s position to be determined with arbitrary precision. Conversely, if one wants to determine position with arbitrary precision, one must introduce a different experimental arrangement, a position-measuring apparatus; but that arrangement necessarily leaves the momentum indeterminate.
%
%The deeper point is that Popper’s attempted reconstruction of a past trajectory illegitimately combines determinations that belong to different measurement contexts. Quantum mechanics does not license the attribution of a single context-independent classical trajectory to the electron. A determinate momentum belongs to a momentum context; a determinate position belongs to a position context. The uncertainty relation is the quantitative expression of the incompatibility between these two possible contexts of determination. It states how much indeterminacy in (q) is necessarily entailed by an arrangement that makes (p) determinate, and conversely how much indeterminacy in (p) is entailed by an arrangement that makes (q) determinate.


%One may reconstruct a quantum history only within a specified context of determination. A sequence of momentum measurements at different times can define a momentum history: the electron may be assigned determinate momenta (p(t_1), p(t_2), p(t_3),\ldots) relative to a momentum-measurement context. Conversely, a sequence of position measurements can define a position history: the electron may be assigned determinate positions (q(t_1), q(t_2), q(t_3),\ldots) relative to a position-measurement context. But there is no common context in which both histories can be combined into a single classical trajectory (q(t),p(t)). Popper’s reconstruction illegitimately treats determinations belonging to distinct and incompatible contexts as if they were fragments of one underlying path. Hermann’s point is not that retrospective reconstruction is impossible as such, but that it is always context-bound: one can reconstruct the causal course appropriate to the actual measurement arrangement, not a context-independent trajectory underlying all possible arrangements.

%She therefore distinguishes two aspects of the wave function. On the one hand, its probabilistic interpretation concerns the transition from one observational context to another. In this respect, it specifies what can be said about possible future observations in a way that corresponds to classical, intuitive description. On the other hand, the wave function also provides the specifically quantum-mechanical description of the state of an individual system relative to the present observational context.

%For this second function, the wave function is not merely a probability distribution over unknown classical particles. It is an exact state-description relative to a fixed observational context. Within such a context, classical double determinations such as position plus momentum are defined only above the limits fixed by the uncertainty relations. If the observation has maximal sharpness relative to that context, nothing remains undetermined in that context. The wave function then gives a complete description of the physical state, independent of probability statements. Hermann supports this by noting that such a description is fixed by specifying the observables and eigenvalues for which the wave function is an eigenstate. Popper’s mistake, then, is to recognize only the probabilistic side of the wave function and to overlook its complementary role as an exact contextual state-description of an individual system. This is why he believes that individual systems may possess more determinate states than the uncertainty relations permit, even though such a belief ignores the quantum-mechanical limits on the applicability of classical concepts.

%Popper's momentum filter presupposes a corpuscular picture: electrons are treated as particles with definite momenta, and the apparatus selects those whose momentum lies within a specified range. Hermann's objection is that this description abstracts from the physical conditions under which such a selection is possible. If the momentum selection is made precise, the electron must be treated as a wave with a correspondingly well-defined wavelength and direction; but such a wave is spatially delocalized. Conversely, any attempt to localize the electron by means of a narrow aperture introduces diffraction and thereby a spread in momentum. Thus the wave aspect is not an additional complication imposed on Popper's experiment from outside; it is already implicit in the very conditions that make the alleged momentum filter a possible measuring arrangement.



%Hermann ends by asking \vW whether he agrees with this account and whether the wave-particle dualism can be experimentally demonstrated not only for large ensembles, but also in processes involving only a few, or even individual, photons or electrons. This question is important for her argument because she wants to show that dualism, and therefore the uncertainty relations, are not merely ensemble-level constraints but belong already to the interpretation of single quantum processes.

Hermann argues that Popper’s appeal to retrospective trajectory calculations rests on a misleading appearance of precision. Such calculations seem to determine the state of a system for a past interval with greater accuracy than the uncertainty relations allow, for example by combining two position measurements, or a position measurement with a momentum measurement. Popper uses these constructions to suggest that the uncertainty relations can be bypassed. Hermann’s diagnosis is that this appearance arises because the calculations ignore the wave-particle dualism. They restrict themselves to the intuitive features of one classical picture, usually the corpuscular picture, and then use two different observations to fix those features. But each individual quantum measurement can be interpreted only by bringing both wave and particle pictures into play. Since no single measurement permits either picture to be applied without restriction, the trajectory calculation illegitimately combines determinations that belong to different observational contexts.

For Hermann, this is precisely what makes the reconstructed trajectory physically meaningless. Quantum mechanics permits state descriptions only relative to a given observational context. A calculation that unifies determinations drawn from different contexts into one description therefore lacks physical significance. This lack of significance shows itself in the fact that, insofar as the reconstruction exceeds the uncertainty relations, it is experimentally uncontrollable and yields no basis for prediction. She contrasts these fictive trajectory calculations with the legitimate inferences required in the interpretation of a measurement. Such inferences are necessary and admissible because they respect the dualism of wave and particle pictures. In the microscope example, the propagation of light through the microscope is described in wave terms, while the collision between light quantum and electron is described in corpuscular terms. The decisive point is that a single observation of the measuring apparatus determines how the two pictures limit one another.

Popper’s arrangement is a momentum-measurement context. It may select electrons with an arbitrarily narrow range of momenta, and in the ideal limit it therefore permits an arbitrarily precise determination of $p_x$. But this very arrangement deprives the electron of a determinate position: a sharply selected momentum corresponds, in wave-mechanical terms, to a plane wave, and a plane wave is spatially delocalized. Conversely, if one wishes to determine the electron’s position with arbitrary precision, one must introduce a different experimental arrangement, one that localizes the electron. But such localization produces a corresponding indeterminacy in momentum. Hermann’s objection to Popper’s reconstruction of past trajectories rests precisely on this point. Popper combines determinations that are meaningful only in different experimental contexts, as though they were partial determinations of one and the same classical trajectory. Hermann denies that this is legitimate. The momentum value belongs to the momentum context, the position value to the position context; they cannot be amalgamated into a single context-independent path of the electron.



%The necessity of choosing between a position-measuring and a momentum-measuring apparatus is not an additional principle over and above the uncertainty relations; it is their experimental expression. Formally, $\hat{x}$ and $\hat{p}_x$ do not commute and therefore have no common eigenstates. Experimentally, this means that no single arrangement can make both $x$ and $p_x$ simultaneously determinate with arbitrary precision. A momentum apparatus prepares or selects a state that approximates an eigenstate of $\widehat{p}_x$, and thereby leaves $x$ indeterminate; a position apparatus prepares or selects a state that approximates an eigenstate of $\widehat{x}$, and thereby leaves $p_x$ indeterminate. Hermann's objection to Popper is that he treats values obtained under these mutually exclusive arrangements as if they could be recombined into one context-indepen trajectory. For Hermann, this recombination is illegitimate: the determinate value is inseparable from the experimental context that makes it determinate.


%Hermann’s response to Popper depends on two connected forms of context-dependence. The first is the context-dependence of measurement: different quantities require different experimental arrangements. A momentum-measuring apparatus can determine p_x with arbitrary precision, but only at the cost of leaving x indeterminate; a position-measuring apparatus can determine x with arbitrary precision, but only at the cost of leaving p_x indeterminate. This prevents Popper from combining a momentum value obtained in one context with a position value obtained in another as if both were partial determinations of a single classical trajectory. The second is the internal structure of each context: within any given arrangement, the uncertainty relation is enforced by the necessary coordination of wave and particle descriptions. A precise momentum selection corresponds to an approximately plane wave and therefore to spatial delocalization; a precise position selection produces diffraction and therefore a spread in momentum. Thus Hermann’s contextualism is not merely the claim that different apparatuses measure different quantities. It is also the claim that each apparatus realizes, in its own way, the wave–particle conditions that make the uncertainty relations unavoidable.

%First, there is inter-context dependence: different quantities require different apparatuses. A momentum measurement and a position measurement are not two readings taken within one and the same experimental arrangement. They require mutually exclusive arrangements

%Second, there is intra-context dependence: within each context, the uncertainty relation is enforced by the coordination of wave and particle descriptions. Even a momentum context is not purely corpuscular. To select a sharp momentum is, in wave terms, to select an approximately plane wave, which makes position indeterminate. Conversely, a position context is not purely spatial or corpuscular. To localize the electron is to introduce wave effects such as diffraction, which makes momentum indeterminate.


%If the observation is made in the image plane, the wave description involves a spherical wave spreading from a point. Correspondingly, the particle description yields a collision with relatively sharp position but blurred momentum exchange. If the observation is made in the focal plane, the wave description involves a parallel bundle of rays. Correspondingly, the particle description yields a sharply determined momentum exchange but a blurred position. Thus the two pictures regulate one another within one observational context. This mutual regulation preserves both the unity of the observational context and the validity of the uncertainty relations. On this basis, Hermann proposes to correct Heisenberg’s remark that it is merely a matter of taste whether one assigns physical reality to such retrospective trajectories. For her, the issue is not taste but contextual legitimacy: these trajectory calculations neglect the essential quantum-mechanical relation between every physical determination and its observational context. For that reason, they become physically meaningless.

%ETP block the argument by denying that the two measurements on particle A can be combined into a single well-defined trajectory description. Once the energy of A is measured, the particle is no longer available as a corpuscle with a determinate spatiotemporal path from which one can reconstruct the relevant arrival time in the required sense. The energy measurement does not simply reveal a preexisting energy while leaving intact the full trajectory information. It changes the conditions under which later time-of-arrival claims can be meaningfully connected to the original emission process.

%Hermann presents Popper’s account as a revealing curiosity. According to Popper, the uncertainty relations rest on the idea that measurement necessarily involves an exchange of energy between the object and the measuring apparatus. This exchange changes the object’s state, so that measurement always reveals a state that has just been destroyed by the very act of measuring. On this view, a measurement cannot tell us the state of the object after the measurement and therefore cannot serve as a basis for prediction.
%
%Hermann treats this as a serious misunderstanding. Popper’s view reduces the uncertainty relations to a crude disturbance thesis: measurement disrupts the object, so the measured state is no longer available afterward. But for Hermann, the issue is not simply that measurement destroys a pre-existing state. The deeper point is that a measurement places the system in a definite observational context, relative to which a new quantum-mechanical state-description is possible.
%
%Her ironic comment is that Popper’s view becomes intelligible only if he takes the fictive retrospective trajectory calculations to be the only complete measurements. In such calculations, one tries to reconstruct a full particle path by combining data from different observations, and the result is indeed useless for prediction because it does not describe the system in any legitimate observational context. Popper then generalizes from this illegitimate case to measurement as such.
%
%Thus Hermann’s criticism is that Popper misunderstands both the uncertainty relations and the function of measurement. A genuine quantum measurement does not merely reveal a state and destroy it; it yields a state-description relative to a new observational context, and within that context it can ground further, though contextually constrained, predictions.


%It meant not the contextual reconstruction of a measurement result, as in Hermann, but the continued possibility that individual particles follow determinate paths, even if \qm, as a statistical theory, does not yet provide the non-statistical laws governing those paths.

This asymmetry in reception is revealing of the changing relation between philosophy and physics. Popper's general philosophy was broad enough to absorb the failure of his interpretation of \qm. His account of conjecture, criticism, and scientific rationality could survive even when his quantum-mechanical arguments came to appear technically fragile and, in retrospect, largely misguided. Hermann's case was the reverse. Her interpretation of \qm was remarkably penetrating, but it was embedded in a Kantian-Friesian framework that did not survive as a compelling general philosophy of science. As a result, her local insight could not easily detach itself from the broader system that carried it, and it was largely forgotten until the recent revival of historically informed philosophy of physics. The subsequent history has partly reversed the terms of comparison. Grand synthetic philosophies of science have become harder to sustain, and Popper's general vision no longer enjoys the authority it once had. At the same time, philosophy of physics has become increasingly technical, historically sensitive, and attentive to the formal and experimental details of quantum theory. In this new context, Hermann's analysis has become newly legible. Popper's philosophy survived because it was larger than his interpretation of \qm; Hermann's interpretation of \qm was forgotten because it was not larger than her philosophy. Yet, as the discipline changed, the very technical precision that once limited her reception became the reason for her rediscovery.


%The asymmetry also has a temporal dimension. Popper's general philosophy was broad enough to absorb the failure of his interpretation of \qm: his account of conjecture, criticism, and scientific rationality could survive even when his quantum-mechanical arguments came to appear technically fragile and largely misguided. Hermann's case was the reverse. Her interpretation of \qm was remarkably penetrating, but it was embedded in a Kantian-Friesian framework that did not survive as a compelling synthetic philosophy of science. Her local insight therefore could not easily detach itself from the broader system that carried it, and it was largely forgotten until the recent revival of historically informed philosophy of physics. In retrospect, this reversal is instructive. Grand synthetic philosophies of science have become harder to sustain, while philosophy of physics has become increasingly technical, historically sensitive, and attentive to the formal and experimental details of quantum theory. In that changed landscape, Hermann's analysis has become newly legible. Popper's philosophy survived because it was larger than his interpretation of \qm; Hermann's interpretation of \qm was rediscovered because philosophy of physics eventually became technical enough to recognize its force.


%This asymmetry in reception is revealing of the changing relation between philosophy and physics. Popper's general philosophy was broad enough to absorb the failure of his interpretation of \qm. His account of conjecture, criticism, and scientific rationality could survive even when his quantum-mechanical arguments came to appear technically fragile and, in retrospect, largely misguided. Hermann's case was the reverse. Her interpretation of \qm was remarkably penetrating, but it was embedded in a Kantian-Friesian framework that did not survive as a compelling general philosophy of science. As a result, her local insight could not easily detach itself from the broader system that carried it, and it was largely forgotten until the recent revival of historically informed philosophy of physics. The subsequent history has partly reversed the terms of comparison. Grand synthetic philosophies of science have become harder to sustain, and Popper's general vision no longer enjoys the authority it once had. At the same time, philosophy of physics has become increasingly technical, historically sensitive, and attentive to the formal and experimental details of quantum theory. In this new context, Hermann's analysis has become newly legible. Popper's philosophy survived because it was larger than his interpretation of \qm; Hermann's interpretation of \qm was forgotten because it was not larger than her philosophy. Yet, as the discipline changed, the very technical precision that once limited her reception became the reason for her rediscovery.


%%This asymmetry in reception is revealing of the complex and evolving relation between philosophy and the sciences, and physics in particular. Popper's technically fragile and ultimately misguided reflections on \qm were embedded in a broader philosophical vision that transformed certain Kantian-Friesian motifs into a historically consequential account of scientific rationality. His failure as an analytic philosopher of physics could therefore be offset by the enormous success of his synthetic philosophy of science. Hermann's mathematically sophisticated interpretation of \qm, by contrast, was embedded in a Kantian-Friesian framework that was just as ambitious as Popper's, but did not survive as a compelling general philosophy of science. Her failure as a synthetic philosopher seems to have carried with it her penetrating and technically informed interpretation of \qm, which was largely forgotten until recently. Indeed, the tables seem to have turned. Grand synthetic philosophies of science have become harder to sustain, and Popper's philosophy has fallen out of fashion. At the same time, philosophy of physics has become increasingly technical and detail-oriented, and Hermann's analysis has acquired new visibility. Yet the tables may turn again, as philosophy of physics risks losing not only the ambition to provide a general philosophy of science, but even the ambition to articulate a philosophical vision of physics as a whole.
%
%%Yet the tables may turn again. If grand synthetic philosophies of science have become increasingly difficult to sustain, the opposite danger now confronts philosophy of physics: its very technical sophistication may lead it to be absorbed into the foundations of physics. In that case, philosophy of physics risks losing not only the ambition to provide a general philosophy of science, but even the ambition to articulate a philosophical vision of physics as a whole. The contrast between Popper and Hermann is instructive in this respect. Popper's achievement reminds us of the power, but also the danger, of an architectonic philosophy that can survive its local scientific failures. Hermann's achievement reminds us of the power, but also the danger, of a technically exact philosophy of physics that may fail to detach its insights from a broader framework, or may later be valued only as a contribution to the foundations of physics. The task, perhaps, is not to choose between synthetic philosophy and analytic philosophy of physics, but to hold the two demands together: philosophical breadth without technical irresponsibility, and technical precision without the loss of philosophical ambition.
%
%
%% whose interpretation of \qm was technically fragile and, in retrospect, largely misguided. As a trained mathematician and a close reader of the quantum formalism, Hermann was the more penetrating \s{analytic} philosopher of physics. Her local insight therefore could not easily detach itself from Hermann's Kantian-Friesian outlook was no less ambitious than Popper's, but did not prove equally durable as a \s{synthetic} philosophy of science. 
%
%%Popper's system survived despite his quantum mechanics; Hermann's quantum mechanics survives despite her system.
%
%%
%%
%%Hermann's interpretation of \qm, by contrast, remains strikingly penetrating even though the philosophical framework in which she embedded it no longer appears equally persuasive. 
%
%%The contrast raises a general question about how philosophers should be judged: by the scope and durability of their architectonic vision, which is howeer increasingly hard to achieve (as Popper's of 
%%
%%or by the precision and accuracy of their local analyses.
%
%
%%Popper's broader philosophy of science could absorb the failure of his interpretation of \qm, because it offered a general and adaptable image of scientific rationality. Hermann's interpretation of \qm, by contrast, remains strikingly penetrating even though the Kantian-Friesian framework in which she embedded it no longer appears equally persuasive.
%
%
%%A powerful philosophical system may illuminate science while misunderstanding one of its central theories; conversely, a historically superseded system may still contain a remarkably precise diagnosis of a concrete scientific problem.
%
%%her attempt and her more general into ethics and politics remained. Popper's philosophical outlook survived despite his interpretation of \qm 
%
%%, of which he was rightly ashamed; Hermann's interpretation of \qm survives despite her philosophical outlook. Hermann, a trained mathematician, was the better philosopher of physics in the modern sense of the expression: she grasped technical aspects of the quantum formalism that Popper failed to understand. Popper, however, turned out to be the more influential philosopher of science: his general account of scientific method proved more adaptable. 
%
%
%The contrast may also be cast as a difference between two philosophical virtues. Popper was, in Schliesser's sense, the more successful synthetic philosopher: he transformed elements of the Kantian-Friesian tradition into a broad, flexible, and historically consequential vision of scientific rationality. Even where his interpretation of \qm failed, it remained embedded in a general philosophy of conjecture, criticism, and realism that could survive the failure of that local application. Hermann's achievement was different. Her Kantian-Friesian system did not prove equally durable as a synthetic philosophy of science. Yet, as a trained mathematician and a close reader of the quantum formalism, she was the more penetrating analytic philosopher of physics: she saw that the central issue was not the recovery of hidden trajectories, but the contextual conditions under which physical quantities can be meaningfully ascribed. Popper's system survived despite his quantum mechanics; Hermann's quantum mechanics survives despite her system. The comparison suggests that philosophers should not be judged only by the durability of their architectonic vision, nor only by the accuracy of their local analyses. A powerful philosophical system may illuminate science while misunderstanding one of its central theories; conversely, a historically superseded system may still contain a remarkably precise diagnosis of a concrete scientific problem.
%
%
%The contrast may be formulated as follows: Popper succeeded as a synthetic philosopher despite his limits as an analytic philosopher of physics; Hermann, conversely, succeeded as an analytic philosopher of physics despite the limits of her synthetic philosophical framework.  The comparison suggests that philosophical achievement should be judged along more than one axis: by the power of a general vision, but also by the precision with which it illuminates a concrete scientific problem.
%
%Popper succeeded as a synthetic philosopher despite his limits as an analytic philosopher of physics; Hermann, conversely, succeeded as an analytic philosopher of physics despite the limits of her synthetic philosophical framework.
%
%
%
%This asymmetry raises a broader question about how philosophers of science should be judged. Popper's general philosophical vision proved extraordinarily durable, even though his specific interpretation of \qm was technically fragile and, in retrospect, largely misguided. Hermann's case was almost the reverse. Her Kantian-Friesian framework did not survive as a compelling general philosophy of science; yet, precisely in her engagement with \qm, she displayed a philosophical penetration that Popper lacked. She understood that the central issue was not simply whether quantum mechanics left room for hidden trajectories, but under what experimental and conceptual conditions physical quantities could be meaningfully ascribed. Popper was therefore the greater philosopher of science in the architectonic sense: he offered a general image of scientific rationality that continued to shape philosophical debate. Hermann was the greater philosopher of physics in the local and diagnostic sense: she grasped the conceptual structure of the quantum formalism with greater accuracy and depth. The contrast suggests that the value of a philosopher cannot be measured only by the survival of a system. Sometimes a failed philosophical framework can contain a remarkably precise insight into a particular science, just as a powerful general philosophy can coexist with a serious misunderstanding of the science it seeks to interpret.
%
%
%%Popper's philosophy can still be appreciated, even by those who dismiss his interpretation of \qm as a historical curiosity; conversely, Hermann's interpretation of quantum mechanics has come to appear be appreciated as penetrating, even to those who do not share her peculiar form of \NKm.
%
%
%%I have reconstructed Popper’s early quantum-mechanical writings, including the unpublished manuscripts and the broader correspondence with Weisskopf, Heisenberg, von Weizsäcker, and Einstein, in Giovanelli [year]. In the present paper I restrict the discussion to the Popper–von Weizsäcker exchange around Popper’s 1934 note, since my aim is to clarify von Weizsäcker’s mediating role between Popper’s trajectory reconstruction and Hermann’s measurement reconstruction.
%
%%In an earlier article, I reconstructed Popper’s 1934 thought experiment and noted that Hermann’s criticism of Popper deserved separate treatment. The present paper takes up that suggestion. It returns to the Popper–von Weizsäcker exchange not in order to repeat the earlier reconstruction, but to clarify von Weizsäcker’s mediating role between Popper’s trajectory reconstruction and Hermann’s measurement reconstruction.
%
%%For a fuller reconstruction of Popper’s 1934 thought experiment and its place in his early interpretation of quantum mechanics, see Giovanelli [year]. The present article returns to this episode only insofar as it bears on von Weizsäcker’s role as a mediating figure between Popper and Hermann. Some technical details of Popper’s arrangement must therefore be recalled, but the focus here is on von Weizsäcker’s response and on its relevance for Hermann’s measurement-theoretical reconstruction.
%
%%Interestingly, both came Friesian-Kantian background, that however they declined in very different way.
%
%%Two ways to save the law of causality.   trajectory recosntrictins ,,, measurement reconstricton
%
%%Popper’s strategy is trajectory reconstruction. He accepts that quantum mechanics does not allow the preparation of ensembles sharper than a pure quantum state, but he insists that this does not exclude the retrospective or indirect determination of the path of an individual particle. His \s{non-prognostic} measurements are meant to show that, after a scattering event, one may reconstruct where the particle was and what momentum it had, and that this reconstruction may even yield testable predictions concerning later detections. Causality is preserved because the individual event can still be inserted into a continuous space-time history.
%%
%%Hermann’s strategy is measurement reconstruction. She does not try to restore a trajectory hidden behind the quantum description. Rather, she argues that once a definite measurement result has occurred, quantum mechanics allows one to reconstruct the causal conditions under which this result arose. The relevant causal nexus is not a context-independent particle path but the specific interaction between object and apparatus. Causality is preserved because each actual result can be explained retrospectively within the experimental context that made it possible. 
%
%%These two different way to save causality: Popper through reconstruction of the particle's path, Hermann, through reconstruction of the measurement context.
%
%%The decisive issue becomes whether there can be elements of physical reality that are not represented in the quantum-mechanical state function. Thus Einstein shifts the debate from the defense of causality to the criterion of completeness.


%If $S$ is placed at the image plane, $S=d_i$, the lens maps object points in $O$ onto image points on $S$, so that a mark on the screen is interpretable as position information. If, instead, $S$ is placed at the focal plane, $S=f$, the lens maps directions of propagation onto points on $S$, so that a mark on the screen is interpretable as momentum information. The distinction is therefore not a difference between two lenses, but a difference between two placements of the screen relative to the same lens.


%image plane selects the position representation; focal plane selects the momentum representation.

%If $S=d_i$, the screen is placed in the image plane and the mark primarily gives position information. However, by reducing the aperture angle $\epsilon$, one narrows the range of possible photon directions and thereby gains more information about the momentum transfer. The price is a loss of positional resolution, since a smaller aperture makes the image of $P$ less sharp.
%
%Conversely, if $S=f$, the screen is placed in the focal plane and the mark primarily gives momentum information, because it selects a definite direction of propagation. However, by increasing the aperture angle $\epsilon$, one admits a wider range of rays from the scattering region and thereby gains more information about the position from which the photon was scattered. The price is a loss of momentum resolution, since a larger aperture corresponds to a wider range of possible photon directions.

%If you choose $S=d_i$, the apparatus is arranged to give position information. A smaller aperture $\epsilon$ does indeed reduce the range of possible photon directions, so in that sense it gives you somewhat better information about the momentum transfer. But the price is that you lose position information, not momentum information.


%- In the Image Plane: The lens converts the angular information (the cone $\epsilon$ ) into a spatial point (the spot $P^{\prime}$ ). The larger the cone $\epsilon$, the smaller the spot.
%- In the Focal Plane: The lens converts the spatial information (the location $P$ ) into angular information (the spread across the focal plane). The smaller the original source $P$ is, the wider the spread across the focal plane.

%In the new inertial state this is represented by $OS$,  Thus the invariant length is not preserved throughout the dynamical process as an ordinary rigid length; rather, it is recovered by the undisturbed rod in its final inertial state.



%A Hermannian answer would be that the wave function characterizes the system only relative to a given observational context. A Bohr-like answer would say that the two arrangements define two different phenomena, not two alternative descriptions of one and the same phenomenon. Thus Einstein sees different descriptions of the same physical state, whereas Bohr and Hermann would see different context-bound phenomena in which the very conditions for state attribution have changed.



%Popper's trajectory is the historical analogue of EPR's simultaneous elements of reality: both result from converting context-bound inferences into context-independent determinations, that one can th $p(t)$. That HEramn that ... and the smae for EPR. Relin once, more clearly, that Popper did not seem to have appreciated Einstein's argument requires not to determine the state of $B$ independently of observation, but on se it does not require the joint possession of $p$ and $q$.ly that the state of $B$ remain the \emph{same}, to determine the state $B$.  

%Yet different measurements on the first system allow different wave functions to be assigned to the second. Thus one and the same real state appears to correspond to different wave functions, and the wave function is therefore incomplete, that is, the wave function cannot represent the real state of the system. That the situation describes, much more eiect. Weizäcker-Hermann \TE.  One might immediately recognize the similarity of this situation to Hermann's treatment of the $\gamma$-ray thought experiment. Hermann did appreciate the very fact that:
%
%$$\psi_{AB}=\psi_A\left(x_1\right) \chi_B\left(x_2\right)$$ 
%
%but is linear combination sum of   
%
%$$\Psi=\sum_{m, n} c_{m n} \psi_m^{(1)}\left(x_1\right) \phi_n^{(2)}\left(x_2\right)$$
%
%Thus precis, the situati  that botehr Einstien but express as temproa, rather the same after collisiton one cie ot measur, and the respticti;  Einstein woild find thi, taht no own that real state of the after the lectorn, has been prestcteley chad. afet. 




%EPR criterion of realyt= certainty of prediction ⇒ reality.
%
%Einstein criterion of samee= same system + no physical disturbance → same state.

%The publication of EPR gave raise; however, the existence of two versions of the argument was not unknown to those who reacted, that EPR version was superogaoty. So the responses were indeed. Of coruse, that would have inflsued at least Heisenberg's reply. In his letter to Bohr on 19 September, Heisenberg and posiibly Bohr appeals to the $\gamma$-ray microscope in order to justify the precise way in which the cut renders impossible the existence of hidden variables. And so Heisenberg, like Hermann, comes to see the microscope thought experiment as more than a mere illustration of the position-momentum uncertainty relation, but as a means of showing the incompletability of the entire theory. Hermann explicitly reaches \the{same system, different wave functions} point. Hermann’s point would be that q and p are not two independent labels that may be attached to the same individual event whenever we can obtain them by clever retrodiction. Each is meaningful only within the experimental context that defines the corresponding quantity. A position-relevant measurement context lets us reconstruct a position; a momentum-relevant context lets us reconstruct a momentum relation. But the two contexts do not jointly define one classical state of the individual electron. Hwoev, that that Einsteun dd nit ry, the same evente two different wave fuctions; hoever, even here the $\gamma$-ray tougth experiuekt  shows, 


%The Bohr-Hermann reply to Einstein is therefore not merely a rejection of the EPR criterion of reality. It is a rejection of the deeper criterion of sameness presupposed by Einstein's argument: the idea that the same undisturbed system must possess the same real state independently of the experimental context. This rejection amounts to an abandonment of what Pauli later called the ideal of the detached observer. Einstein's realism requires a state of the system independent of observation; Hermann's contextualism makes state attribution dependent on the conditions under which the system is experimentally accessible.
%
%Einstein grants, for the sake of argument, that no definite $q$ or $p$ need be attributed to the distant system. His point is only that the system’s real state cannot depend on which measurement is later chosen elsewhere. The contextualist reply is that this already assumes what quantum mechanics calls into question: that an individual subsystem has a context-independent real state. For Hermann, Heisenberg, and Bohr, the different wave functions do not reveal the incompleteness of $\psi$; they reveal the illegitimacy of abstracting a subsystem-state from the whole observational context in which alone state attribution has physical meaning.
%
%This means not that the electron has no state at all, but that it does not have its own separable complete state. It has only a reduced mixed state, while the pure state belongs to the correlated electron–photon system. We can save relativistic no-signaling and deny any mechanical disturbance, but we must abandon the classical idea that spatially separated subsystems always possess independent complete states.
%
%Physics, for Einstein, requires the possibility of decomposing the world into relatively independent systems, each possessing its own real state. Otherwise, the notion of a physical system becomes unstable. The orthodox reply distinguishes no-action from separability. It preserves the former by pointing to the invariance of marginal statistics, but it abandons the latter because the individual subsystem has no complete state of its own. Einstein’s incompleteness argument then reappears as the claim that a theory which only secures statistical closure, without describing the individual separated system, is not a complete theory of physical reality. For the orthodox quantum theorist, the charge of “incompleteness” is not necessarily damaging, because the theory was never meant to describe an individual system in Einstein’s sense. It gives the complete rules for possible predictions, including all marginal and conditional statistics. The orthodox reply is satisfied with statistical closure: the marginal state of the distant system is unchanged, so there is no action at a distance. Einstein accepts this formal point but denies that it solves the foundational problem. If quantum mechanics cannot assign a complete real state to the individual separated system, then it is incomplete. And if it is incomplete, one should seek a deeper theory that restores an objective description of individual systems.
%
%The distinction between proper and improper mixtures is internal to the quantum formalism. It tells us how the mixed state is generated within the theory. But it does not, by itself, settle whether the subsystem possesses an underlying real state not represented by the quantum state. Therefore it cannot serve as a non-question-begging refutation of Einstein’s incompleteness claim. The appeal to improper mixtures sharpens the orthodox reply but does not defeat Einstein. It shows that, within quantum mechanics, the local mixed state of a subsystem is not a proper ignorance mixture over quantum states. But Einstein can reinterpret precisely this fact as a symptom of incompleteness: the theory supplies only the reduced statistical state, not the real state of the individual subsystem. Thus the debate cannot be settled merely by invoking improper mixtures; it turns on whether the quantum state is assumed to be complete. You have not shown that the individual system has no real state. You have only stipulated that the quantum state assigned to the subsystem is not to be understood as ignorance about such a state.
%
%In Pauli’s words, the observer’s \s{subjective element} is the free choice among exclusive arrangements. Once the arrangement is chosen, the individual outcome is no longer under the observer’s control. That is where objectivity enters: not as deterministic predictability, but as the uncontrollability of the individual result under a fixed arrangement.
%
%The \s{orthodox} side says: once thie experimental arrangement is fixed, the state is fixed; if the arrangement changes, the state attribution changes. There is no further question about a context-independent state of the individual system behind these attributions. The \s{objectivity} of the physicist’s world, as far as indeterminate individual events are concerned, results from the assumption, apparently well supported empirically, that the physicist cannot influence these events—that is, the indeterminate individual result of his measurements, once he has chosen his experimental arrangement.
% 
%%The statistics do not arise merely because the possible arrangements are incompatible. They arise because, once a state has been defined by one arrangement, a later incompatible arrangement introduces an uncontrollable interaction with the system. Since this interaction cannot be specified in detail, quantum mechanics renounces deterministic prediction of individual outcomes and replaces it with statistical predictions for ensembles.
% 
%One can not even speak about the same state. Einstein says: that is exactly the problem. If the \s{state} depends on the arrangement, then it is not the real state of the system. It is a description of our mode of access, or of the statistical predictions licensed by a context. Therefore quantum mechanics is incomplete as a description of individual physical reality. 
% 



%The public EPR argument rests on a criterion of reality: if a quantity can be predicted with certainty without physically disturbing the system, then there is an element of reality corresponding to that quantity. Applied to an entangled pair, this criterion appears to license the attribution of both momentum and position to the distant system, depending on which measurement is performed on its partner. If the corresponding wave functions are eigenstates of the momentum and position operators, $\psi_p$ and $\psi_q$, then the public EPR argument seems to restore precisely the simultaneous sharp values excluded by the uncertainty relations.
%
%Einstein’s own argument is weaker. It does not require the joint attribution of sharp values of non-commuting observables. It rests instead on a criterion of sameness: if the same system is not physically disturbed, then its real state must remain the same. Thus, if different measurements on one system allow different wave functions, say $\psi$ and $\psi’$, to be assigned to the distant system, while that distant system is not physically disturbed, then one and the same real state appears to correspond to different wave-function descriptions. Einstein’s conclusion is therefore not that the uncertainty relations are violated, but that the wave function cannot be a complete description of the real state.
%
%Hermann’s discussion of Weizsäcker's $\gamma$-ray microscope was written before Einstein’s own EPR argument. Yet might have influece kind of response that Heisenberg, and possibly Bohr, would later give to EPR. The central move is to deny that the \s{absence of disturbance} suffices for context-independent state attribution. Against the public EPR argument, this blocks the inference from certainty of prediction without disturbance to the \emph{reality} of a corresponding physical property. This of course blocks also Popper's attempt to to into a single trajectory.  
%
%Against Einstein’s private argument, it blocks the weaker inference from numerical identity of the system and absence of disturbance to \emph{sameness} of real state. The state, for Hermann, is not fixed by the isolated object alone, but by the experimental arrangement that determines which predicates can be meaningfully applied. Even if there is no mechanical disturbance, the choice of measurement can still alter the conditions under which predictions about the distant system are meaningful. That is the \s{non-mechanical} disturbance: not a physical kick, but a change in the experimental conditions defining what can be said.
%
%Indeed, that same state of the electron of the collision is described by different wave functions.
%It is probably here a clear philosophical disagreement: Einstein is bothered by the possibility of one real state corresponding to several wave functions. Hermann is not, because she rejects the assumption that the relevant state is a context-independent real state of the isolated system. For her, the wave function belongs to a specified observational context; different contexts can therefore yield different wave functions for the same system without implying incompleteness. This is probably the conclusion that Einstein wanted, that it cannot depend on the choice of the measuring apparatus. That Einstein found that it depends on experimental arrangement, or on the way one looks at it. That has to state that the world would accept the same for macroscopic systems; that just, also exists. In his later years, he insisted on this point, that the world would accept the same for macroscopic systems. 
%
%After the particle has passed the slit, one can decide which kind of experimental arrangement is to be used for its description: one that licenses a position attribution, or one that licenses a momentum attribution. But one cannot combine the two resulting descriptions into a single state of the same individual particle, because the conditions under which they are meaningful are mutually exclusive.
%
%\text{state} = \text{state relative to an experimental arrangement}. \text{state} = \text{real condition of an individual system, independent of how it is observed}.

%Bohr’s example concerns the attribution of a state to the particle at the time of its passage through the slit. But this attribution is not independent of the later-described experimental arrangement: a fixed diaphragm licenses a position description, while a movable diaphragm licenses a momentum description. Hence the alternative descriptions cannot be combined into one classical state of the particle before the slit.

%Hermann’s electron-photon case and Bohr’s slit case both concern the retrospective description of an individual process. In each case, one may reconstruct either a position-relevant or a momentum-relevant description of the particle’s past state, but only relative to an observational context that excludes the other reconstruction.


 

%Bohm: restore individual systems, but accept nonlocality  Everett: remove collapse and universalize the wave function


%%Bohr notion of non-mechanical disturbance 
%
%%No physical disturbance is not sufficient, by itself, to secure either a context-independent element of reality or a context-independent identity of state. 
%
%
%%
%%Here the restoposes that, depending on the observational context, one can obtain different wave functions for the same system and the same instant, namely the electron immediately after collision with the light quantum. They add that, for Hermann, the quantum-mechanical characterization does not pertain to the physical systemin itself, independently of the observation by which one obtains knowledge of it. 
%
%
%
%
%%Hermann’s position is strikingly close in spirit to the later consistent-histories approach: both deny that one may freely combine descriptions belonging to mutually incompatible experimental or inferential contexts. A quantum system cannot be assigned, without qualification, all the properties that would be available in different classical descriptions. Rather, meaningful ascriptions of properties are valid only within a determinate framework.
%
%%For Griffiths, the key point is the single-framework rule. One may reason within a consistent family of histories, but one may not combine propositions drawn from incompatible frameworks. Thus, for example, one framework may describe a particle as having a well-defined position history, another as having a well-defined momentum history, but one cannot fuse these into a single trajectory with both sharp position and sharp momentum at every relevant time.

%Hermann can be read as offering a philosophical ancestor of the consistent-histories attitude, though not of its formal machinery. Like Griffiths, she blocks the illicit construction of a single global history from elements belonging to incompatible descriptions. Unlike Griffiths, however, she frames the point in terms of the conditions under which causal explanation remains possible: the causal chain can be reconstructed only relative to the observational context that fixes the relevant state description.




%
%The public EPR argument relies on a criterion of reality: certainty of prediction, together with the absence of physical disturbance, licenses the attribution of a corresponding element of reality.
%
%Einstein’s deeper argument relies instead on a criterion of sameness: if it is the same system and there is no physical disturbance, then the real state of that system must remain the same.
%
%The Bohr-Hermann-Weizsäcker reply challenges both inferences. 


%This is not an operational violation of the uncertainty relation, since one does not jointly measure position and momentum on $B$ itself. It is rather a realist challenge to the completeness of the wave function: the real state of $B$ appears to contain determinate features that the quantum-mechanical description cannot simultaneously express.
%Einstein's second argument is much more vulnerable to a Hermann-style criticism. It depends on combining two counterfactual experimental possibilities: if one measures $A$ in one way, one can predict the position of $B$; if one measures $A$ in another way, one can predict the momentum of $B$. Einstein then infers that, since either prediction can be made without disturbing $B$, both quantities must belong to the real physical state of $B$. Hermann would reject precisely this inference. For her, each prediction is valid only relative to the observational context that makes it possible. A position measurement licenses a position attribution, while excluding the context in which momentum could be sharply determined; a momentum measurement licenses a momentum attribution, while excluding the context in which position could be sharply determined. Thus one cannot combine the two counterfactual attributions into a single context-independent state of $B$.


%EInstien explaint, that after, say, a collision, two particles are described by the joint, nonfactorizable wave function\footnote{As is well known, the term \s{entanglement} was first introduced by Schrödinger (\cite{Schroedinger1935} and \cite{Schrodinger1935b}) in the ensuing months} $\psi(x_1, x_2)$, with sharp relative position ($x_B - x_A = 0$) and sharp total momentum ($p_A + p_B = 0$)---but spread individual particle momenta.  It is undeniable that the sitatuin resameb, Popper's argumet, and even more clear Herman's \TE. As we that 
%
%If one measures, say, the momentum $p_A$ of particle $A$, its state collapses into the corresponding momentum eigenstate $\psi_A$; quantum mechanics then assigns to subsystem $B$ a conditional momentum eigenstate $\psi_B$. For this state, however, the position $x_B$ is completely undetermined, which blocks any attempt to beat the \UR, as Popper had hoped. Einstein's point, one ede, to emausre postion  taht rather, was that the specific form of $\psi_B$ would have been different had one performed a position measurement on system $A$. Hoever, th system $B$ cannot have chafed. IS that, that measurem on $A$ would real state of $B$ This means that two different $\psi$-functions correspond to one and the same physical state of $B$, which itself has not changed physically; thys cannot be of a comete \q{It is therefore not possible to regard the $\psi$-function as the complete description of the state of the system.}\letteraep{Einstein}{Popper}{11}{9}{1935}[19-130].
%
%%That that cannot be the same system, that it cannot be that same state, cannot change the that measureing cannot chane the satte of electron in the past. Hower, that version of Hoeer, the the publsied version makes as step fowd. 
%
%However, since Popper had asked him to summarize the published EPR paper, Einstein also reconstruct that argumet, Once one can prect would meaomnt say $p$ one  $-p$, the smae for position; both position and momentum of $B$ can be predicted with certainty, according to the infamous EPR \s{criterion of reality}, \emph{both} quantities must correspond to something that exists in physical reality. This situation, because of the \IR{}, finds no counterpart in the formalism of \qm{}.\letteraep{Einstein}{Popper}{11}{9}{1935}[19-130] Although the argument was again meant to convey the incompleteness of the theory, it also seems to conflict with the spirit of the \IR{}. Indeed, Popper\autocite[244\fn{4*}]{Popper1959} would later describe the EPR argument, as presented in this letter, as a weaker (no prediction for both position and momentum simultaneously) but correct (both position and momentum are simultaneously part of reality) version of his own failed ETP-like \TE{}. That criteiro  of reality.
%
%
%
%%that simle simlaritiew iwht Einstien's argumet, at least a collsion. The goal of the argument was to show that \qm leads to the consequence that location and momentum must be ‘elements of reality’ against the uncertainty relations.  The argument did not provide a prediction, as he hoped, but that they are both real for a single system against the \UR. 
%
%In the public EPR argument, the question is whether a clever arrangement permits the attribution of both $p$ and $q$ values to the same individual system. In Popper's argument, the question is whether successive measurements and retrospective inference permit the attribution of both $p$ and $q$ values to the same individual history. Hermann's answer to both is structurally similar: quantum mechanics does not license the fusion of mutually exclusive context-bound $\mathrm{d}^{\star} \sim$ rminations into a single classical description.
%
%Hermann rejects that inference. For her, the validity of an attribution depends not only on predictive or retrodictive certainty, but on the experimental context that makes the relevant predicate applicable.
%
%$p$ is predictable with certainty and $q$ is predictable with certainty therefore, EPR conclude, both correspond to elements of reality of the same distant system. For Popper, the target is retrospective trajectory reconstruction: $q\left(t_0\right)$ can be inferred retrospectively and $p\left(t_0\right)$ can be inferred retrospectively therefore, Popper concludes, the particle has a determinate classical history. Of corue, also the make  predctuon 